\documentclass[11pt]{article}
    \newcommand{\myname}{Ryan Hughes}
    \newcommand{\githuburl}{https://github.com/rhughes74/Assignments} % link to assignment02.ipynb
    \newif\ifshowcode
    \showcodetrue   % set to \showcodefalse to drop the code appendix

    \usepackage{fancyhdr}
    \setlength{\headheight}{14pt}   % avoids a "headheight too small" warning
    \pagestyle{fancy}
    \fancyhf{}                      % clear the default header/footer
    \lhead{\myname}                 % left
    \chead{Assignment 2}            % center
    \rhead{CNET 5051, Fall 2026}    % right
    \cfoot{\thepage}                % page number, centered at the bottom
    \renewcommand{\headrulewidth}{0.4pt}   % thin line under the header
    \usepackage[breakable]{tcolorbox}
    \usepackage{parskip} % Stop auto-indenting (to mimic markdown behaviour)
    
    \usepackage[breakable]{tcolorbox}   % breakable lets a box split across pages

    \newtcolorbox{problemstmt}{
    breakable,
    colback=gray!12,       % fill color
    colframe=gray!50,      % border color
    boxrule=0.5pt,
    arc=2pt,
    left=6pt, right=6pt, top=4pt, bottom=4pt
    }

    % Basic figure setup, for now with no caption control since it's done
    % automatically by Pandoc (which extracts ![](path) syntax from Markdown).
    \usepackage{graphicx}
    % Keep aspect ratio if custom image width or height is specified
    \setkeys{Gin}{keepaspectratio}
    % Maintain compatibility with old templates. Remove in nbconvert 6.0
    \let\Oldincludegraphics\includegraphics
    % Ensure that by default, figures have no caption (until we provide a
    % proper Figure object with a Caption API and a way to capture that
    % in the conversion process - todo).
    \usepackage{caption}
    \DeclareCaptionFormat{nocaption}{}
    \captionsetup{format=nocaption,aboveskip=0pt,belowskip=0pt}

    \usepackage{float}
    \floatplacement{figure}{H} % forces figures to be placed at the correct location
    \usepackage{xcolor} % Allow colors to be defined
    \usepackage{enumerate} % Needed for markdown enumerations to work
    \usepackage{geometry} % Used to adjust the document margins
    \usepackage{amsmath} % Equations
    \usepackage{amssymb} % Equations
    \usepackage{textcomp} % defines textquotesingle
    % Hack from http://tex.stackexchange.com/a/47451/13684:
    \AtBeginDocument{%
        \def\PYZsq{\textquotesingle}% Upright quotes in Pygmentized code
    }
    \usepackage{upquote} % Upright quotes for verbatim code
    \usepackage{eurosym} % defines \euro

    \usepackage{iftex}
    \ifPDFTeX
        \usepackage[T1]{fontenc}
        \IfFileExists{alphabeta.sty}{
              \usepackage{alphabeta}
          }{
              \usepackage[mathletters]{ucs}
              \usepackage[utf8x]{inputenc}
          }
    \else
        \usepackage{fontspec}
        \usepackage{unicode-math}
    \fi

    \usepackage{fancyvrb} % verbatim replacement that allows latex
    \usepackage{grffile} % extends the file name processing of package graphics
                         % to support a larger range
    \makeatletter % fix for old versions of grffile with XeLaTeX
    \@ifpackagelater{grffile}{2019/11/01}
    {
      % Do nothing on new versions
    }
    {
      \def\Gread@@xetex#1{%
        \IfFileExists{"\Gin@base".bb}%
        {\Gread@eps{\Gin@base.bb}}%
        {\Gread@@xetex@aux#1}%
      }
    }
    \makeatother
    \usepackage[Export]{adjustbox} % Used to constrain images to a maximum size
    \adjustboxset{max size={0.9\linewidth}{0.9\paperheight}}

    % The hyperref package gives us a pdf with properly built
    % internal navigation ('pdf bookmarks' for the table of contents,
    % internal cross-reference links, web links for URLs, etc.)
    \usepackage{hyperref}
    % The default LaTeX title has an obnoxious amount of whitespace. By default,
    % titling removes some of it. It also provides customization options.
    \usepackage{titling}
    \usepackage{longtable} % longtable support required by pandoc >1.10
    \usepackage{booktabs}  % table support for pandoc > 1.12.2
    \usepackage{array}     % table support for pandoc >= 2.11.3
    \usepackage{calc}      % table minipage width calculation for pandoc >= 2.11.1
    \usepackage[inline]{enumitem} % IRkernel/repr support (it uses the enumerate* environment)
    \usepackage[normalem]{ulem} % ulem is needed to support strikethroughs (\sout)
                                % normalem makes italics be italics, not underlines
    \usepackage{soul}      % strikethrough (\st) support for pandoc >= 3.0.0
    \usepackage{mathrsfs}
    

    
    % Colors for the hyperref package
    \definecolor{urlcolor}{rgb}{0,.145,.698}
    \definecolor{linkcolor}{rgb}{.71,0.21,0.01}
    \definecolor{citecolor}{rgb}{.12,.54,.11}

    % ANSI colors
    \definecolor{ansi-black}{HTML}{3E424D}
    \definecolor{ansi-black-intense}{HTML}{282C36}
    \definecolor{ansi-red}{HTML}{E75C58}
    \definecolor{ansi-red-intense}{HTML}{B22B31}
    \definecolor{ansi-green}{HTML}{00A250}
    \definecolor{ansi-green-intense}{HTML}{007427}
    \definecolor{ansi-yellow}{HTML}{DDB62B}
    \definecolor{ansi-yellow-intense}{HTML}{B27D12}
    \definecolor{ansi-blue}{HTML}{208FFB}
    \definecolor{ansi-blue-intense}{HTML}{0065CA}
    \definecolor{ansi-magenta}{HTML}{D160C4}
    \definecolor{ansi-magenta-intense}{HTML}{A03196}
    \definecolor{ansi-cyan}{HTML}{60C6C8}
    \definecolor{ansi-cyan-intense}{HTML}{258F8F}
    \definecolor{ansi-white}{HTML}{C5C1B4}
    \definecolor{ansi-white-intense}{HTML}{A1A6B2}
    \definecolor{ansi-default-inverse-fg}{HTML}{FFFFFF}
    \definecolor{ansi-default-inverse-bg}{HTML}{000000}

    % common color for the border for error outputs.
    \definecolor{outerrorbackground}{HTML}{FFDFDF}

    % commands and environments needed by pandoc snippets
    % extracted from the output of `pandoc -s`
    \providecommand{\tightlist}{%
      \setlength{\itemsep}{0pt}\setlength{\parskip}{0pt}}
    \DefineVerbatimEnvironment{Highlighting}{Verbatim}{commandchars=\\\{\}}
    % Add ',fontsize=\small' for more characters per line
    \newenvironment{Shaded}{}{}
    \newcommand{\KeywordTok}[1]{\textcolor[rgb]{0.00,0.44,0.13}{\textbf{{#1}}}}
    \newcommand{\DataTypeTok}[1]{\textcolor[rgb]{0.56,0.13,0.00}{{#1}}}
    \newcommand{\DecValTok}[1]{\textcolor[rgb]{0.25,0.63,0.44}{{#1}}}
    \newcommand{\BaseNTok}[1]{\textcolor[rgb]{0.25,0.63,0.44}{{#1}}}
    \newcommand{\FloatTok}[1]{\textcolor[rgb]{0.25,0.63,0.44}{{#1}}}
    \newcommand{\CharTok}[1]{\textcolor[rgb]{0.25,0.44,0.63}{{#1}}}
    \newcommand{\StringTok}[1]{\textcolor[rgb]{0.25,0.44,0.63}{{#1}}}
    \newcommand{\CommentTok}[1]{\textcolor[rgb]{0.38,0.63,0.69}{\textit{{#1}}}}
    \newcommand{\OtherTok}[1]{\textcolor[rgb]{0.00,0.44,0.13}{{#1}}}
    \newcommand{\AlertTok}[1]{\textcolor[rgb]{1.00,0.00,0.00}{\textbf{{#1}}}}
    \newcommand{\FunctionTok}[1]{\textcolor[rgb]{0.02,0.16,0.49}{{#1}}}
    \newcommand{\RegionMarkerTok}[1]{{#1}}
    \newcommand{\ErrorTok}[1]{\textcolor[rgb]{1.00,0.00,0.00}{\textbf{{#1}}}}
    \newcommand{\NormalTok}[1]{{#1}}

    % Additional commands for more recent versions of Pandoc
    \newcommand{\ConstantTok}[1]{\textcolor[rgb]{0.53,0.00,0.00}{{#1}}}
    \newcommand{\SpecialCharTok}[1]{\textcolor[rgb]{0.25,0.44,0.63}{{#1}}}
    \newcommand{\VerbatimStringTok}[1]{\textcolor[rgb]{0.25,0.44,0.63}{{#1}}}
    \newcommand{\SpecialStringTok}[1]{\textcolor[rgb]{0.73,0.40,0.53}{{#1}}}
    \newcommand{\ImportTok}[1]{{#1}}
    \newcommand{\DocumentationTok}[1]{\textcolor[rgb]{0.73,0.13,0.13}{\textit{{#1}}}}
    \newcommand{\AnnotationTok}[1]{\textcolor[rgb]{0.38,0.63,0.69}{\textbf{\textit{{#1}}}}}
    \newcommand{\CommentVarTok}[1]{\textcolor[rgb]{0.38,0.63,0.69}{\textbf{\textit{{#1}}}}}
    \newcommand{\VariableTok}[1]{\textcolor[rgb]{0.10,0.09,0.49}{{#1}}}
    \newcommand{\ControlFlowTok}[1]{\textcolor[rgb]{0.00,0.44,0.13}{\textbf{{#1}}}}
    \newcommand{\OperatorTok}[1]{\textcolor[rgb]{0.40,0.40,0.40}{{#1}}}
    \newcommand{\BuiltInTok}[1]{{#1}}
    \newcommand{\ExtensionTok}[1]{{#1}}
    \newcommand{\PreprocessorTok}[1]{\textcolor[rgb]{0.74,0.48,0.00}{{#1}}}
    \newcommand{\AttributeTok}[1]{\textcolor[rgb]{0.49,0.56,0.16}{{#1}}}
    \newcommand{\InformationTok}[1]{\textcolor[rgb]{0.38,0.63,0.69}{\textbf{\textit{{#1}}}}}
    \newcommand{\WarningTok}[1]{\textcolor[rgb]{0.38,0.63,0.69}{\textbf{\textit{{#1}}}}}
    \makeatletter
    \newsavebox\pandoc@box
    \newcommand*\pandocbounded[1]{%
      \sbox\pandoc@box{#1}%
      % scaling factors for width and height
      \Gscale@div\@tempa\textheight{\dimexpr\ht\pandoc@box+\dp\pandoc@box\relax}%
      \Gscale@div\@tempb\linewidth{\wd\pandoc@box}%
      % select the smaller of both
      \ifdim\@tempb\p@<\@tempa\p@
        \let\@tempa\@tempb
      \fi
      % scaling accordingly (\@tempa < 1)
      \ifdim\@tempa\p@<\p@
        \scalebox{\@tempa}{\usebox\pandoc@box}%
      % scaling not needed, use as it is
      \else
        \usebox{\pandoc@box}%
      \fi
    }
    \makeatother

    % Define a nice break command that doesn't care if a line doesn't already
    % exist.
    \def\br{\hspace*{\fill} \\* }
    % Math Jax compatibility definitions
    \def\gt{>}
    \def\lt{<}
    \let\Oldtex\TeX
    \let\Oldlatex\LaTeX
    \renewcommand{\TeX}{\textrm{\Oldtex}}
    \renewcommand{\LaTeX}{\textrm{\Oldlatex}}
    % Document parameters
    % Document title
    \title{Hughes}
    
    
    
    
    
    
    
% Pygments definitions
\makeatletter
\def\PY@reset{\let\PY@it=\relax \let\PY@bf=\relax%
    \let\PY@ul=\relax \let\PY@tc=\relax%
    \let\PY@bc=\relax \let\PY@ff=\relax}
\def\PY@tok#1{\csname PY@tok@#1\endcsname}
\def\PY@toks#1+{\ifx\relax#1\empty\else%
    \PY@tok{#1}\expandafter\PY@toks\fi}
\def\PY@do#1{\PY@bc{\PY@tc{\PY@ul{%
    \PY@it{\PY@bf{\PY@ff{#1}}}}}}}
\def\PY#1#2{\PY@reset\PY@toks#1+\relax+\PY@do{#2}}

\@namedef{PY@tok@w}{\def\PY@tc##1{\textcolor[rgb]{0.73,0.73,0.73}{##1}}}
\@namedef{PY@tok@c}{\let\PY@it=\textit\def\PY@tc##1{\textcolor[rgb]{0.24,0.48,0.48}{##1}}}
\@namedef{PY@tok@cp}{\def\PY@tc##1{\textcolor[rgb]{0.61,0.40,0.00}{##1}}}
\@namedef{PY@tok@k}{\let\PY@bf=\textbf\def\PY@tc##1{\textcolor[rgb]{0.00,0.50,0.00}{##1}}}
\@namedef{PY@tok@kp}{\def\PY@tc##1{\textcolor[rgb]{0.00,0.50,0.00}{##1}}}
\@namedef{PY@tok@kt}{\def\PY@tc##1{\textcolor[rgb]{0.69,0.00,0.25}{##1}}}
\@namedef{PY@tok@o}{\def\PY@tc##1{\textcolor[rgb]{0.40,0.40,0.40}{##1}}}
\@namedef{PY@tok@ow}{\let\PY@bf=\textbf\def\PY@tc##1{\textcolor[rgb]{0.67,0.13,1.00}{##1}}}
\@namedef{PY@tok@nb}{\def\PY@tc##1{\textcolor[rgb]{0.00,0.50,0.00}{##1}}}
\@namedef{PY@tok@nf}{\def\PY@tc##1{\textcolor[rgb]{0.00,0.00,1.00}{##1}}}
\@namedef{PY@tok@nc}{\let\PY@bf=\textbf\def\PY@tc##1{\textcolor[rgb]{0.00,0.00,1.00}{##1}}}
\@namedef{PY@tok@nn}{\let\PY@bf=\textbf\def\PY@tc##1{\textcolor[rgb]{0.00,0.00,1.00}{##1}}}
\@namedef{PY@tok@ne}{\let\PY@bf=\textbf\def\PY@tc##1{\textcolor[rgb]{0.80,0.25,0.22}{##1}}}
\@namedef{PY@tok@nv}{\def\PY@tc##1{\textcolor[rgb]{0.10,0.09,0.49}{##1}}}
\@namedef{PY@tok@no}{\def\PY@tc##1{\textcolor[rgb]{0.53,0.00,0.00}{##1}}}
\@namedef{PY@tok@nl}{\def\PY@tc##1{\textcolor[rgb]{0.46,0.46,0.00}{##1}}}
\@namedef{PY@tok@ni}{\let\PY@bf=\textbf\def\PY@tc##1{\textcolor[rgb]{0.44,0.44,0.44}{##1}}}
\@namedef{PY@tok@na}{\def\PY@tc##1{\textcolor[rgb]{0.41,0.47,0.13}{##1}}}
\@namedef{PY@tok@nt}{\let\PY@bf=\textbf\def\PY@tc##1{\textcolor[rgb]{0.00,0.50,0.00}{##1}}}
\@namedef{PY@tok@nd}{\def\PY@tc##1{\textcolor[rgb]{0.67,0.13,1.00}{##1}}}
\@namedef{PY@tok@s}{\def\PY@tc##1{\textcolor[rgb]{0.73,0.13,0.13}{##1}}}
\@namedef{PY@tok@sd}{\let\PY@it=\textit\def\PY@tc##1{\textcolor[rgb]{0.73,0.13,0.13}{##1}}}
\@namedef{PY@tok@si}{\let\PY@bf=\textbf\def\PY@tc##1{\textcolor[rgb]{0.64,0.35,0.47}{##1}}}
\@namedef{PY@tok@se}{\let\PY@bf=\textbf\def\PY@tc##1{\textcolor[rgb]{0.67,0.36,0.12}{##1}}}
\@namedef{PY@tok@sr}{\def\PY@tc##1{\textcolor[rgb]{0.64,0.35,0.47}{##1}}}
\@namedef{PY@tok@ss}{\def\PY@tc##1{\textcolor[rgb]{0.10,0.09,0.49}{##1}}}
\@namedef{PY@tok@sx}{\def\PY@tc##1{\textcolor[rgb]{0.00,0.50,0.00}{##1}}}
\@namedef{PY@tok@m}{\def\PY@tc##1{\textcolor[rgb]{0.40,0.40,0.40}{##1}}}
\@namedef{PY@tok@gh}{\let\PY@bf=\textbf\def\PY@tc##1{\textcolor[rgb]{0.00,0.00,0.50}{##1}}}
\@namedef{PY@tok@gu}{\let\PY@bf=\textbf\def\PY@tc##1{\textcolor[rgb]{0.50,0.00,0.50}{##1}}}
\@namedef{PY@tok@gd}{\def\PY@tc##1{\textcolor[rgb]{0.63,0.00,0.00}{##1}}}
\@namedef{PY@tok@gi}{\def\PY@tc##1{\textcolor[rgb]{0.00,0.52,0.00}{##1}}}
\@namedef{PY@tok@gr}{\def\PY@tc##1{\textcolor[rgb]{0.89,0.00,0.00}{##1}}}
\@namedef{PY@tok@ge}{\let\PY@it=\textit}
\@namedef{PY@tok@gs}{\let\PY@bf=\textbf}
\@namedef{PY@tok@ges}{\let\PY@bf=\textbf\let\PY@it=\textit}
\@namedef{PY@tok@gp}{\let\PY@bf=\textbf\def\PY@tc##1{\textcolor[rgb]{0.00,0.00,0.50}{##1}}}
\@namedef{PY@tok@go}{\def\PY@tc##1{\textcolor[rgb]{0.44,0.44,0.44}{##1}}}
\@namedef{PY@tok@gt}{\def\PY@tc##1{\textcolor[rgb]{0.00,0.27,0.87}{##1}}}
\@namedef{PY@tok@err}{\def\PY@bc##1{{\setlength{\fboxsep}{\string -\fboxrule}\fcolorbox[rgb]{1.00,0.00,0.00}{1,1,1}{\strut ##1}}}}
\@namedef{PY@tok@kc}{\let\PY@bf=\textbf\def\PY@tc##1{\textcolor[rgb]{0.00,0.50,0.00}{##1}}}
\@namedef{PY@tok@kd}{\let\PY@bf=\textbf\def\PY@tc##1{\textcolor[rgb]{0.00,0.50,0.00}{##1}}}
\@namedef{PY@tok@kn}{\let\PY@bf=\textbf\def\PY@tc##1{\textcolor[rgb]{0.00,0.50,0.00}{##1}}}
\@namedef{PY@tok@kr}{\let\PY@bf=\textbf\def\PY@tc##1{\textcolor[rgb]{0.00,0.50,0.00}{##1}}}
\@namedef{PY@tok@bp}{\def\PY@tc##1{\textcolor[rgb]{0.00,0.50,0.00}{##1}}}
\@namedef{PY@tok@fm}{\def\PY@tc##1{\textcolor[rgb]{0.00,0.00,1.00}{##1}}}
\@namedef{PY@tok@vc}{\def\PY@tc##1{\textcolor[rgb]{0.10,0.09,0.49}{##1}}}
\@namedef{PY@tok@vg}{\def\PY@tc##1{\textcolor[rgb]{0.10,0.09,0.49}{##1}}}
\@namedef{PY@tok@vi}{\def\PY@tc##1{\textcolor[rgb]{0.10,0.09,0.49}{##1}}}
\@namedef{PY@tok@vm}{\def\PY@tc##1{\textcolor[rgb]{0.10,0.09,0.49}{##1}}}
\@namedef{PY@tok@sa}{\def\PY@tc##1{\textcolor[rgb]{0.73,0.13,0.13}{##1}}}
\@namedef{PY@tok@sb}{\def\PY@tc##1{\textcolor[rgb]{0.73,0.13,0.13}{##1}}}
\@namedef{PY@tok@sc}{\def\PY@tc##1{\textcolor[rgb]{0.73,0.13,0.13}{##1}}}
\@namedef{PY@tok@dl}{\def\PY@tc##1{\textcolor[rgb]{0.73,0.13,0.13}{##1}}}
\@namedef{PY@tok@s2}{\def\PY@tc##1{\textcolor[rgb]{0.73,0.13,0.13}{##1}}}
\@namedef{PY@tok@sh}{\def\PY@tc##1{\textcolor[rgb]{0.73,0.13,0.13}{##1}}}
\@namedef{PY@tok@s1}{\def\PY@tc##1{\textcolor[rgb]{0.73,0.13,0.13}{##1}}}
\@namedef{PY@tok@mb}{\def\PY@tc##1{\textcolor[rgb]{0.40,0.40,0.40}{##1}}}
\@namedef{PY@tok@mf}{\def\PY@tc##1{\textcolor[rgb]{0.40,0.40,0.40}{##1}}}
\@namedef{PY@tok@mh}{\def\PY@tc##1{\textcolor[rgb]{0.40,0.40,0.40}{##1}}}
\@namedef{PY@tok@mi}{\def\PY@tc##1{\textcolor[rgb]{0.40,0.40,0.40}{##1}}}
\@namedef{PY@tok@il}{\def\PY@tc##1{\textcolor[rgb]{0.40,0.40,0.40}{##1}}}
\@namedef{PY@tok@mo}{\def\PY@tc##1{\textcolor[rgb]{0.40,0.40,0.40}{##1}}}
\@namedef{PY@tok@ch}{\let\PY@it=\textit\def\PY@tc##1{\textcolor[rgb]{0.24,0.48,0.48}{##1}}}
\@namedef{PY@tok@cm}{\let\PY@it=\textit\def\PY@tc##1{\textcolor[rgb]{0.24,0.48,0.48}{##1}}}
\@namedef{PY@tok@cpf}{\let\PY@it=\textit\def\PY@tc##1{\textcolor[rgb]{0.24,0.48,0.48}{##1}}}
\@namedef{PY@tok@c1}{\let\PY@it=\textit\def\PY@tc##1{\textcolor[rgb]{0.24,0.48,0.48}{##1}}}
\@namedef{PY@tok@cs}{\let\PY@it=\textit\def\PY@tc##1{\textcolor[rgb]{0.24,0.48,0.48}{##1}}}

\def\PYZbs{\char`\\}
\def\PYZus{\char`\_}
\def\PYZob{\char`\{}
\def\PYZcb{\char`\}}
\def\PYZca{\char`\^}
\def\PYZam{\char`\&}
\def\PYZlt{\char`\<}
\def\PYZgt{\char`\>}
\def\PYZsh{\char`\#}
\def\PYZpc{\char`\%}
\def\PYZdl{\char`\$}
\def\PYZhy{\char`\-}
\def\PYZsq{\char`\'}
\def\PYZdq{\char`\"}
\def\PYZti{\char`\~}
% for compatibility with earlier versions
\def\PYZat{@}
\def\PYZlb{[}
\def\PYZrb{]}
\makeatother


    % For linebreaks inside Verbatim environment from package fancyvrb.
    \makeatletter
        \newbox\Wrappedcontinuationbox
        \newbox\Wrappedvisiblespacebox
        \newcommand*\Wrappedvisiblespace {\textcolor{red}{\textvisiblespace}}
        \newcommand*\Wrappedcontinuationsymbol {\textcolor{red}{\llap{\tiny$\m@th\hookrightarrow$}}}
        \newcommand*\Wrappedcontinuationindent {3ex }
        \newcommand*\Wrappedafterbreak {\kern\Wrappedcontinuationindent\copy\Wrappedcontinuationbox}
        % Take advantage of the already applied Pygments mark-up to insert
        % potential linebreaks for TeX processing.
        %        {, <, #, %, $, ' and ": go to next line.
        %        _, }, ^, &, >, - and ~: stay at end of broken line.
        % Use of \textquotesingle for straight quote.
        \newcommand*\Wrappedbreaksatspecials {%
            \def\PYGZus{\discretionary{\char`\_}{\Wrappedafterbreak}{\char`\_}}%
            \def\PYGZob{\discretionary{}{\Wrappedafterbreak\char`\{}{\char`\{}}%
            \def\PYGZcb{\discretionary{\char`\}}{\Wrappedafterbreak}{\char`\}}}%
            \def\PYGZca{\discretionary{\char`\^}{\Wrappedafterbreak}{\char`\^}}%
            \def\PYGZam{\discretionary{\char`\&}{\Wrappedafterbreak}{\char`\&}}%
            \def\PYGZlt{\discretionary{}{\Wrappedafterbreak\char`\<}{\char`\<}}%
            \def\PYGZgt{\discretionary{\char`\>}{\Wrappedafterbreak}{\char`\>}}%
            \def\PYGZsh{\discretionary{}{\Wrappedafterbreak\char`\#}{\char`\#}}%
            \def\PYGZpc{\discretionary{}{\Wrappedafterbreak\char`\%}{\char`\%}}%
            \def\PYGZdl{\discretionary{}{\Wrappedafterbreak\char`\$}{\char`\$}}%
            \def\PYGZhy{\discretionary{\char`\-}{\Wrappedafterbreak}{\char`\-}}%
            \def\PYGZsq{\discretionary{}{\Wrappedafterbreak\textquotesingle}{\textquotesingle}}%
            \def\PYGZdq{\discretionary{}{\Wrappedafterbreak\char`\"}{\char`\"}}%
            \def\PYGZti{\discretionary{\char`\~}{\Wrappedafterbreak}{\char`\~}}%
        }
        % Some characters . , ; ? ! / are not pygmentized.
        % This macro makes them "active" and they will insert potential linebreaks
        \newcommand*\Wrappedbreaksatpunct {%
            \lccode`\~`\.\lowercase{\def~}{\discretionary{\hbox{\char`\.}}{\Wrappedafterbreak}{\hbox{\char`\.}}}%
            \lccode`\~`\,\lowercase{\def~}{\discretionary{\hbox{\char`\,}}{\Wrappedafterbreak}{\hbox{\char`\,}}}%
            \lccode`\~`\;\lowercase{\def~}{\discretionary{\hbox{\char`\;}}{\Wrappedafterbreak}{\hbox{\char`\;}}}%
            \lccode`\~`\:\lowercase{\def~}{\discretionary{\hbox{\char`\:}}{\Wrappedafterbreak}{\hbox{\char`\:}}}%
            \lccode`\~`\?\lowercase{\def~}{\discretionary{\hbox{\char`\?}}{\Wrappedafterbreak}{\hbox{\char`\?}}}%
            \lccode`\~`\!\lowercase{\def~}{\discretionary{\hbox{\char`\!}}{\Wrappedafterbreak}{\hbox{\char`\!}}}%
            \lccode`\~`\/\lowercase{\def~}{\discretionary{\hbox{\char`\/}}{\Wrappedafterbreak}{\hbox{\char`\/}}}%
            \catcode`\.\active
            \catcode`\,\active
            \catcode`\;\active
            \catcode`\:\active
            \catcode`\?\active
            \catcode`\!\active
            \catcode`\/\active
            \lccode`\~`\~
        }
    \makeatother

    \let\OriginalVerbatim=\Verbatim
    \makeatletter
    \renewcommand{\Verbatim}[1][1]{%
        %\parskip\z@skip
        \sbox\Wrappedcontinuationbox {\Wrappedcontinuationsymbol}%
        \sbox\Wrappedvisiblespacebox {\FV@SetupFont\Wrappedvisiblespace}%
        \def\FancyVerbFormatLine ##1{\hsize\linewidth
            \vtop{\raggedright\hyphenpenalty\z@\exhyphenpenalty\z@
                \doublehyphendemerits\z@\finalhyphendemerits\z@
                \strut ##1\strut}%
        }%
        % If the linebreak is at a space, the latter will be displayed as visible
        % space at end of first line, and a continuation symbol starts next line.
        % Stretch/shrink are however usually zero for typewriter font.
        \def\FV@Space {%
            \nobreak\hskip\z@ plus\fontdimen3\font minus\fontdimen4\font
            \discretionary{\copy\Wrappedvisiblespacebox}{\Wrappedafterbreak}
            {\kern\fontdimen2\font}%
        }%

        % Allow breaks at special characters using \PYG... macros.
        \Wrappedbreaksatspecials
        % Breaks at punctuation characters . , ; ? ! and / need catcode=\active
        \OriginalVerbatim[#1,codes*=\Wrappedbreaksatpunct]%
    }
    \makeatother

    % Exact colors from NB
    \definecolor{incolor}{HTML}{303F9F}
    \definecolor{outcolor}{HTML}{D84315}
    \definecolor{cellborder}{HTML}{CFCFCF}
    \definecolor{cellbackground}{HTML}{F7F7F7}

    % prompt
    \makeatletter
    \newcommand{\boxspacing}{\kern\kvtcb@left@rule\kern\kvtcb@boxsep}
    \makeatother
    \newcommand{\prompt}[4]{
        {\ttfamily\llap{{\color{#2}[#3]:\hspace{3pt}#4}}\vspace{-\baselineskip}}
    }
    

    
    % Prevent overflowing lines due to hard-to-break entities
    \sloppy
    % Setup hyperref package
    \hypersetup{
      breaklinks=true,  % so long urls are correctly broken across lines
      colorlinks=true,
      urlcolor=urlcolor,
      linkcolor=linkcolor,
      citecolor=citecolor,
      }
    % Slightly bigger margins than the latex defaults
    
    \geometry{verbose,tmargin=1in,bmargin=1in,lmargin=1in,rmargin=1in}
    
    

\begin{document}
    
    \maketitle
    
    



\textbf{Overall Instructions:} Due by Friday, Oct.~30th, 11:00pm EST.
Please submit a .pdf of your answers via Canvas, including a url to your
.ipynb notebook in a public GitHub repository. Use the provided LATEX
template for your submission. Please name your Canvas submission as
Lastname.pdf, (e.g.~Trujillo.pdf). If you need help with LATEX or
GitHub, please visit during office hours or post questions to the
Canvas. Include your code in a (well-documented, commented, clear)
assignment02.ipynb file. Scores for the code-based questions below will
be based on your written responses and code.
\begin{problemstmt}
\textbf{Problem 1 (20 points)} In this question we will work with
degree-controlled random graph null models.

\begin{enumerate}
\def\labelenumi{(\alph{enumi})}
\item
  The Configuration Model (5 points). First, implement the configuration
  model in Python. Your function should take a list of integers as a
  degree sequence, and return a graph where each vertex has the assigned
  degree. You graph may contain duplicate edges and self-loops.
\item
  The Chung Lu Model (5 points). Next, implement the Chung Lu random
  graph model. Remember that this model creates an edge between nodes i
  and j with independent probabilities defined by vertex degree and the
  total number of edges in the graph: pij = kikj 2m , where 2m := ∑︁ i
  ki. Your implementation should take a list of integers as a degree
  sequence, and return a simple graph with approximately the desired
  degree sequence.
\item
  Model Differences (5 points). For the degree sequence
  degrees={[}8,5,5,4,3,2,2,2,1,1,1,1,1,1,1{]}, generate a thousand ran-
  dom graphs using the Configuration Model and the Chung Lu model, then:
\item
  Plot the average adjacency matrix for each model. This should be a
  heatmap of: ⟨A⟩ = 1 N N∑︂ i=1 Ai where N is the number of generated
  random graphs, and Ai is the adjacency matrix of the ith graph. (Hint:
  matplotlib's imshow is useful for displaying matrix data.)
\end{enumerate}

\begin{enumerate}
\def\labelenumi{(\roman{enumi})}
\setcounter{enumi}{1}
\item
  Choose two graph metrics: one that should be approximately the same
  across graphs gen- erated by both models, and one that should diverge
  significantly between the two models. For each metric, make a plot
  (histogram, kernel-density estimate, or cumulative distribution
  function) showing the distribution of the metric across all generated
  graphs from each model. Note: For this stage you may remove self-loops
  and duplicate edges from your configuration model graphs if it helps
  make a direct comparison to Chung Lu graphs.
\item
  Discussion (5 points). Discuss the above results briefly. Why do the
  models produce similar results under some metrics and different
  results under others?
\end{enumerate}
\end{problemstmt}
    \begin{tcolorbox}[breakable, size=fbox, boxrule=1pt, pad at break*=1mm,colback=cellbackground, colframe=cellborder]
\prompt{In}{incolor}{1}{\boxspacing}
\begin{Verbatim}[commandchars=\\\{\}]
\PY{k+kn}{import}\PY{+w}{ }\PY{n+nn}{networkx}\PY{+w}{ }\PY{k}{as}\PY{+w}{ }\PY{n+nn}{nx}
\PY{k+kn}{import}\PY{+w}{ }\PY{n+nn}{matplotlib}\PY{n+nn}{.}\PY{n+nn}{pyplot}\PY{+w}{ }\PY{k}{as}\PY{+w}{ }\PY{n+nn}{plt}
\PY{k+kn}{import}\PY{+w}{ }\PY{n+nn}{numpy}\PY{+w}{ }\PY{k}{as}\PY{+w}{ }\PY{n+nn}{np}
\PY{k+kn}{from}\PY{+w}{ }\PY{n+nn}{collections}\PY{+w}{ }\PY{k+kn}{import} \PY{n}{Counter}
\PY{k+kn}{from}\PY{+w}{ }\PY{n+nn}{itertools}\PY{+w}{ }\PY{k+kn}{import} \PY{n}{combinations}
\PY{k+kn}{import}\PY{+w}{ }\PY{n+nn}{logging}
\PY{k+kn}{import}\PY{+w}{ }\PY{n+nn}{networkx}\PY{+w}{ }\PY{k}{as}\PY{+w}{ }\PY{n+nn}{nx}
\PY{k+kn}{import}\PY{+w}{ }\PY{n+nn}{matplotlib}\PY{n+nn}{.}\PY{n+nn}{pyplot}\PY{+w}{ }\PY{k}{as}\PY{+w}{ }\PY{n+nn}{plt}
\PY{k+kn}{from}\PY{+w}{ }\PY{n+nn}{scipy}\PY{n+nn}{.}\PY{n+nn}{stats}\PY{+w}{ }\PY{k+kn}{import} \PY{n}{binom}
\PY{k+kn}{from}\PY{+w}{ }\PY{n+nn}{scipy}\PY{n+nn}{.}\PY{n+nn}{special}\PY{+w}{ }\PY{k+kn}{import} \PY{n}{comb}

\PY{k}{def}\PY{+w}{ }\PY{n+nf}{configuration\PYZus{}model\PYZus{}from\PYZus{}degree\PYZus{}sequence}\PY{p}{(}\PY{n}{degree\PYZus{}sequence}\PY{p}{,} \PY{n}{return\PYZus{}simple}\PY{o}{=}\PY{k+kc}{True}\PY{p}{)}\PY{p}{:}
\PY{+w}{    }\PY{l+s+sd}{\PYZdq{}\PYZdq{}\PYZdq{}}
\PY{l+s+sd}{    Generate a random graph using the configuration model from a given degree sequence}
\PY{l+s+sd}{    without using the NetworkX built\PYZhy{}in function.}

\PY{l+s+sd}{    The configuration model generates a random graph that preserves the degree }
\PY{l+s+sd}{    sequence of nodes by assigning \PYZdq{}stubs\PYZdq{} or \PYZdq{}half\PYZhy{}edges\PYZdq{} to each node and }
\PY{l+s+sd}{    randomly pairing these stubs to form edges. This process can result in }
\PY{l+s+sd}{    graphs with self\PYZhy{}loops and parallel edges, which can be removed if needed.}

\PY{l+s+sd}{    Parameters}
\PY{l+s+sd}{    \PYZhy{}\PYZhy{}\PYZhy{}\PYZhy{}\PYZhy{}\PYZhy{}\PYZhy{}\PYZhy{}\PYZhy{}\PYZhy{}}
\PY{l+s+sd}{    degree\PYZus{}sequence : list of int}
\PY{l+s+sd}{        A list representing the degree of each node in the graph. The sum of }
\PY{l+s+sd}{        the degrees in this sequence must be even for the configuration model }
\PY{l+s+sd}{        to create a valid graph.}

\PY{l+s+sd}{    Returns}
\PY{l+s+sd}{    \PYZhy{}\PYZhy{}\PYZhy{}\PYZhy{}\PYZhy{}\PYZhy{}\PYZhy{}}
\PY{l+s+sd}{    G : networkx.MultiGraph}
\PY{l+s+sd}{        A multigraph generated from the given degree sequence. The graph may }
\PY{l+s+sd}{        contain self\PYZhy{}loops and parallel edges, which are allowed in the }
\PY{l+s+sd}{        configuration model.}

\PY{l+s+sd}{    Notes}
\PY{l+s+sd}{    \PYZhy{}\PYZhy{}\PYZhy{}\PYZhy{}\PYZhy{}}
\PY{l+s+sd}{    \PYZhy{} This method works by assigning \PYZdq{}stubs\PYZdq{} or \PYZdq{}half\PYZhy{}edges\PYZdq{} to nodes based on }
\PY{l+s+sd}{      their degree and then randomly pairing them to form edges. The resulting }
\PY{l+s+sd}{      graph can have self\PYZhy{}loops and parallel edges.}
\PY{l+s+sd}{    \PYZhy{} Self\PYZhy{}loops and parallel edges can be removed post\PYZhy{}generation if a simple }
\PY{l+s+sd}{      graph is required using NetworkX\PYZsq{}s `nx.Graph(G)`.}
\PY{l+s+sd}{    \PYZhy{} The degree sequence must have an even sum for a valid graph construction. }
\PY{l+s+sd}{      If the sum of the degrees is odd, no graph can be constructed.}

\PY{l+s+sd}{    Time Complexity}
\PY{l+s+sd}{    \PYZhy{}\PYZhy{}\PYZhy{}\PYZhy{}\PYZhy{}\PYZhy{}\PYZhy{}\PYZhy{}\PYZhy{}\PYZhy{}\PYZhy{}\PYZhy{}\PYZhy{}\PYZhy{}\PYZhy{}}
\PY{l+s+sd}{    The time complexity is O(E), where E is the number of edges in the graph.}

\PY{l+s+sd}{    Example}
\PY{l+s+sd}{    \PYZhy{}\PYZhy{}\PYZhy{}\PYZhy{}\PYZhy{}\PYZhy{}\PYZhy{}}
\PY{l+s+sd}{    \PYZgt{}\PYZgt{}\PYZgt{} degree\PYZus{}sequence = [3, 3, 2, 2, 1, 1]}
\PY{l+s+sd}{    \PYZgt{}\PYZgt{}\PYZgt{} G = configuration\PYZus{}model\PYZus{}from\PYZus{}degree\PYZus{}sequence(degree\PYZus{}sequence)}
\PY{l+s+sd}{    \PYZgt{}\PYZgt{}\PYZgt{} nx.is\PYZus{}graphical(degree\PYZus{}sequence)}
\PY{l+s+sd}{    True}
\PY{l+s+sd}{    \PYZdq{}\PYZdq{}\PYZdq{}}

    \PY{c+c1}{\PYZsh{} Check if the degree sequence is valid (sum of degrees must be even)}
    \PY{k}{if} \PY{n+nb}{sum}\PY{p}{(}\PY{n}{degree\PYZus{}sequence}\PY{p}{)} \PY{o}{\PYZpc{}} \PY{l+m+mi}{2} \PY{o}{!=} \PY{l+m+mi}{0}\PY{p}{:}
        \PY{k}{raise} \PY{n+ne}{ValueError}\PY{p}{(}\PY{l+s+s2}{\PYZdq{}}\PY{l+s+s2}{The sum of the degree sequence must be even.}\PY{l+s+s2}{\PYZdq{}}\PY{p}{)}

    \PY{c+c1}{\PYZsh{} Create stubs list: node i appears degree\PYZus{}sequence[i] times}
    \PY{n}{stubs} \PY{o}{=} \PY{p}{[}\PY{p}{]}
    \PY{k}{for} \PY{n}{node}\PY{p}{,} \PY{n}{degree} \PY{o+ow}{in} \PY{n+nb}{enumerate}\PY{p}{(}\PY{n}{degree\PYZus{}sequence}\PY{p}{)}\PY{p}{:}
        \PY{n}{stubs}\PY{o}{.}\PY{n}{extend}\PY{p}{(}\PY{p}{[}\PY{n}{node}\PY{p}{]} \PY{o}{*} \PY{n}{degree}\PY{p}{)}

    \PY{c+c1}{\PYZsh{} Shuffle stubs to randomize the pairing process}
    \PY{n}{np}\PY{o}{.}\PY{n}{random}\PY{o}{.}\PY{n}{shuffle}\PY{p}{(}\PY{n}{stubs}\PY{p}{)}

    \PY{c+c1}{\PYZsh{} Initialize an empty multigraph}
    \PY{n}{G} \PY{o}{=} \PY{n}{nx}\PY{o}{.}\PY{n}{MultiGraph}\PY{p}{(}\PY{p}{)}

    \PY{c+c1}{\PYZsh{} Add nodes to the graph}
    \PY{n}{G}\PY{o}{.}\PY{n}{add\PYZus{}nodes\PYZus{}from}\PY{p}{(}\PY{n+nb}{range}\PY{p}{(}\PY{n+nb}{len}\PY{p}{(}\PY{n}{degree\PYZus{}sequence}\PY{p}{)}\PY{p}{)}\PY{p}{)}

    \PY{c+c1}{\PYZsh{} Pair stubs to create edges}
    \PY{k}{while} \PY{n}{stubs}\PY{p}{:}
        \PY{n}{u} \PY{o}{=} \PY{n}{stubs}\PY{o}{.}\PY{n}{pop}\PY{p}{(}\PY{p}{)}
        \PY{n}{v} \PY{o}{=} \PY{n}{stubs}\PY{o}{.}\PY{n}{pop}\PY{p}{(}\PY{p}{)}

        \PY{c+c1}{\PYZsh{} Add the edge to the graph}
        \PY{n}{G}\PY{o}{.}\PY{n}{add\PYZus{}edge}\PY{p}{(}\PY{n}{u}\PY{p}{,} \PY{n}{v}\PY{p}{)}

    \PY{k}{if} \PY{n}{return\PYZus{}simple}\PY{p}{:}
        \PY{c+c1}{\PYZsh{} Convert the multigraph to a simple graph (remove parallel edges)}
        \PY{n}{G\PYZus{}simple} \PY{o}{=} \PY{n}{nx}\PY{o}{.}\PY{n}{Graph}\PY{p}{(}\PY{n}{G}\PY{p}{)}  \PY{c+c1}{\PYZsh{} This removes parallel edges by default}

        \PY{k}{return} \PY{n}{G\PYZus{}simple}

    \PY{k}{else}\PY{p}{:}
        \PY{k}{return} \PY{n}{G}


\PY{k}{def}\PY{+w}{ }\PY{n+nf}{chung\PYZus{}lu\PYZus{}model\PYZus{}from\PYZus{}degree\PYZus{}sequence}\PY{p}{(}\PY{n}{degree\PYZus{}sequence}\PY{p}{,} \PY{n}{return\PYZus{}simple}\PY{o}{=}\PY{k+kc}{True}\PY{p}{)}\PY{p}{:}
\PY{+w}{    }\PY{l+s+sd}{\PYZdq{}\PYZdq{}\PYZdq{}}
\PY{l+s+sd}{    Generate a random graph using the Chung Lu model from a given degree sequence}
\PY{l+s+sd}{    without using the NetworkX built\PYZhy{}in function.}

\PY{l+s+sd}{    The Chung Lu model generates a random graph that preserves the degree }
\PY{l+s+sd}{    sequence of nodes by...}

\PY{l+s+sd}{    Parameters}
\PY{l+s+sd}{    \PYZhy{}\PYZhy{}\PYZhy{}\PYZhy{}\PYZhy{}\PYZhy{}\PYZhy{}\PYZhy{}\PYZhy{}\PYZhy{}}
\PY{l+s+sd}{    degree\PYZus{}sequence : list of int}
\PY{l+s+sd}{        A list representing the degree of each node in the graph. The sum of }
\PY{l+s+sd}{        the degrees in this sequence must be even for the Chung Lu model }
\PY{l+s+sd}{        to create a valid graph.}

\PY{l+s+sd}{    Returns}
\PY{l+s+sd}{    \PYZhy{}\PYZhy{}\PYZhy{}\PYZhy{}\PYZhy{}\PYZhy{}\PYZhy{}}
\PY{l+s+sd}{    G : networkx.Graph}
\PY{l+s+sd}{        A graph generated from the given degree sequence. }
\PY{l+s+sd}{        }
\PY{l+s+sd}{    Time Complexity}
\PY{l+s+sd}{    \PYZhy{}\PYZhy{}\PYZhy{}\PYZhy{}\PYZhy{}\PYZhy{}\PYZhy{}\PYZhy{}\PYZhy{}\PYZhy{}\PYZhy{}\PYZhy{}\PYZhy{}\PYZhy{}\PYZhy{}}
\PY{l+s+sd}{    The time complexity is O(M+N), where M is the number of edges in the graph and N the number of nodes}
\PY{l+s+sd}{    \PYZdq{}\PYZdq{}\PYZdq{}}
    \PY{c+c1}{\PYZsh{}Instatiate G:}
    \PY{n}{G} \PY{o}{=} \PY{n}{nx}\PY{o}{.}\PY{n}{Graph}\PY{p}{(}\PY{p}{)}
    \PY{n}{N}\PY{o}{=}\PY{n+nb}{len}\PY{p}{(}\PY{n}{degree\PYZus{}sequence}\PY{p}{)} \PY{c+c1}{\PYZsh{}Total number of nodes by taking the length of the degree sequence}
    \PY{n}{G}\PY{o}{.}\PY{n}{add\PYZus{}nodes\PYZus{}from}\PY{p}{(}\PY{n+nb}{range}\PY{p}{(}\PY{n}{N}\PY{p}{)}\PY{p}{)}

    \PY{c+c1}{\PYZsh{} Check if the degree sequence is valid (sum of degrees must be even)}
    \PY{k}{if} \PY{n+nb}{sum}\PY{p}{(}\PY{n}{degree\PYZus{}sequence}\PY{p}{)} \PY{o}{\PYZpc{}} \PY{l+m+mi}{2} \PY{o}{!=} \PY{l+m+mi}{0}\PY{p}{:}
        \PY{k}{raise} \PY{n+ne}{ValueError}\PY{p}{(}\PY{l+s+s2}{\PYZdq{}}\PY{l+s+s2}{The sum of the degree sequence must be even.}\PY{l+s+s2}{\PYZdq{}}\PY{p}{)}

    \PY{c+c1}{\PYZsh{}First calculate the total number of edges in the graph by taking the sum of the degree sequence}
    \PY{n}{total\PYZus{}edges} \PY{o}{=} \PY{n+nb}{sum}\PY{p}{(}\PY{n}{degree\PYZus{}sequence}\PY{p}{)}
    
    \PY{c+c1}{\PYZsh{}Ensure degree sequence is sorted:}
    \PY{n}{degree\PYZus{}sequence}\PY{o}{=}\PY{n}{np}\PY{o}{.}\PY{n}{sort}\PY{p}{(}\PY{n}{degree\PYZus{}sequence}\PY{p}{,}\PY{n}{descending}\PY{o}{=}\PY{k+kc}{True}\PY{p}{)}

    \PY{k}{for} \PY{n}{i} \PY{o+ow}{in} \PY{n+nb}{range}\PY{p}{(}\PY{l+m+mi}{0}\PY{p}{,}\PY{n}{N}\PY{o}{\PYZhy{}}\PY{l+m+mi}{1}\PY{p}{)}\PY{p}{:}\PY{c+c1}{\PYZsh{}for all nodes}
        \PY{n}{j}\PY{o}{=}\PY{n}{i}\PY{o}{+}\PY{l+m+mi}{1}\PY{c+c1}{\PYZsh{}Look at the next possible node to connect to...}
        \PY{c+c1}{\PYZsh{}and, get the probability of connecting to this node:}
        \PY{n}{pij} \PY{o}{=} \PY{n+nb}{min}\PY{p}{(}\PY{p}{(}\PY{n}{degree\PYZus{}sequence}\PY{p}{[}\PY{n}{i}\PY{p}{]} \PY{o}{*} \PY{n}{degree\PYZus{}sequence}\PY{p}{[}\PY{n}{j}\PY{p}{]}\PY{p}{)} \PY{o}{/} \PY{p}{(}\PY{n}{total\PYZus{}edges}\PY{p}{)}\PY{p}{,}\PY{l+m+mi}{1}\PY{p}{)}
        
        \PY{c+c1}{\PYZsh{}while there are additional valid other nodes available to connect to...}
        \PY{k}{while} \PY{n}{j} \PY{o}{\PYZlt{}} \PY{n}{N} \PY{o+ow}{and} \PY{n}{pij} \PY{o}{\PYZgt{}} \PY{l+m+mi}{0}\PY{p}{:} 
            \PY{k}{if} \PY{p}{(}\PY{n}{pij} \PY{o}{\PYZlt{}} \PY{l+m+mi}{1}\PY{p}{)}\PY{p}{:}\PY{c+c1}{\PYZsh{}(i.e., while the probability is not 100\PYZpc{} or 0\PYZpc{})}
                \PY{n}{r} \PY{o}{=} \PY{n}{np}\PY{o}{.}\PY{n}{random}\PY{o}{.}\PY{n}{rand}\PY{p}{(}\PY{p}{)}\PY{c+c1}{\PYZsh{}draw a random chance}
                \PY{c+c1}{\PYZsh{}and calculate the following to grab one of these additional other nodes:}
                \PY{n}{j} \PY{o}{+}\PY{o}{=} \PY{n+nb}{int}\PY{p}{(}\PY{n}{np}\PY{o}{.}\PY{n}{floor}\PY{p}{(}\PY{n}{np}\PY{o}{.}\PY{n}{log}\PY{p}{(}\PY{n}{r}\PY{p}{)}\PY{o}{/}\PY{n}{np}\PY{o}{.}\PY{n}{log}\PY{p}{(}\PY{l+m+mi}{1}\PY{o}{\PYZhy{}}\PY{n}{pij}\PY{p}{)}\PY{p}{)}\PY{p}{)}

            \PY{k}{if} \PY{n}{j} \PY{o}{\PYZlt{}} \PY{n}{N}\PY{p}{:}\PY{c+c1}{\PYZsh{}while there are additional other nodes to process...}
                \PY{c+c1}{\PYZsh{}get the additional other node\PYZsq{}s probability of connecting to:}
                \PY{n}{q} \PY{o}{=} \PY{n+nb}{min}\PY{p}{(}\PY{p}{(}\PY{n}{degree\PYZus{}sequence}\PY{p}{[}\PY{n}{i}\PY{p}{]} \PY{o}{*} \PY{n}{degree\PYZus{}sequence}\PY{p}{[}\PY{n}{j}\PY{p}{]}\PY{p}{)} \PY{o}{/} \PY{p}{(}\PY{n}{total\PYZus{}edges}\PY{p}{)}\PY{p}{,}\PY{l+m+mi}{1}\PY{p}{)}
                \PY{n}{r} \PY{o}{=} \PY{n}{np}\PY{o}{.}\PY{n}{random}\PY{o}{.}\PY{n}{rand}\PY{p}{(}\PY{p}{)}\PY{c+c1}{\PYZsh{}draw a random chance}
                \PY{k}{if} \PY{n}{r} \PY{o}{\PYZlt{}} \PY{n}{q}\PY{o}{/}\PY{n}{pij}\PY{p}{:}\PY{c+c1}{\PYZsh{}if the random chance is less than the probability of connecting to this new other node relative to the simple next in the sequence:}
                    \PY{n}{G}\PY{o}{.}\PY{n}{add\PYZus{}edge}\PY{p}{(}\PY{n}{i}\PY{p}{,}\PY{n}{j}\PY{p}{)}\PY{c+c1}{\PYZsh{}add an edge to the new other node}
                \PY{n}{pij}\PY{o}{=}\PY{n}{q}\PY{c+c1}{\PYZsh{}update the probability of connection with the last processed node}
                \PY{n}{j}\PY{o}{+}\PY{o}{=}\PY{l+m+mi}{1}\PY{c+c1}{\PYZsh{}move forward relative to the new other node}
    \PY{k}{return} \PY{n}{G}

\PY{c+c1}{\PYZsh{}(b) Testing}
\PY{n}{degrees}\PY{o}{=}\PY{p}{[}\PY{l+m+mi}{8}\PY{p}{,}\PY{l+m+mi}{5}\PY{p}{,}\PY{l+m+mi}{5}\PY{p}{,}\PY{l+m+mi}{4}\PY{p}{,}\PY{l+m+mi}{3}\PY{p}{,}\PY{l+m+mi}{2}\PY{p}{,}\PY{l+m+mi}{2}\PY{p}{,}\PY{l+m+mi}{2}\PY{p}{,}\PY{l+m+mi}{1}\PY{p}{,}\PY{l+m+mi}{1}\PY{p}{,}\PY{l+m+mi}{1}\PY{p}{,}\PY{l+m+mi}{1}\PY{p}{,}\PY{l+m+mi}{1}\PY{p}{,}\PY{l+m+mi}{1}\PY{p}{,}\PY{l+m+mi}{1}\PY{p}{]}
\PY{n}{chung\PYZus{}lu\PYZus{}graph} \PY{o}{=} \PY{n}{chung\PYZus{}lu\PYZus{}model\PYZus{}from\PYZus{}degree\PYZus{}sequence}\PY{p}{(}\PY{n}{degrees}\PY{p}{)}
\PY{n+nb}{print}\PY{p}{(}\PY{n}{chung\PYZus{}lu\PYZus{}graph}\PY{o}{.}\PY{n}{degree}\PY{p}{)}
\PY{c+c1}{\PYZsh{}(c) Testing}
\PY{n}{chung\PYZus{}lus} \PY{o}{=} \PY{l+m+mi}{0}
\PY{n}{configs} \PY{o}{=} \PY{l+m+mi}{0}
\PY{n}{num\PYZus{}edges\PYZus{}list\PYZus{}config}\PY{o}{=}\PY{p}{[}\PY{p}{]}
\PY{n}{num\PYZus{}edges\PYZus{}list\PYZus{}chunglu}\PY{o}{=}\PY{p}{[}\PY{p}{]}
\PY{n}{cc\PYZus{}list\PYZus{}config}\PY{o}{=}\PY{p}{[}\PY{p}{]}
\PY{n}{cc\PYZus{}list\PYZus{}chunglu}\PY{o}{=}\PY{p}{[}\PY{p}{]}
\PY{n}{pos} \PY{o}{=} \PY{p}{[}\PY{p}{]}
\PY{k}{for} \PY{n}{i} \PY{o+ow}{in} \PY{n+nb}{range}\PY{p}{(}\PY{l+m+mi}{1}\PY{p}{,}\PY{l+m+mi}{1001}\PY{p}{)}\PY{p}{:}
    \PY{n}{cc\PYZus{}list\PYZus{}config}\PY{o}{.}\PY{n}{append}\PY{p}{(}\PY{n}{nx}\PY{o}{.}\PY{n}{average\PYZus{}clustering}\PY{p}{(}\PY{p}{(}\PY{n}{configuration\PYZus{}model\PYZus{}from\PYZus{}degree\PYZus{}sequence}\PY{p}{(}\PY{n}{degrees}\PY{p}{)}\PY{p}{)}\PY{p}{)}\PY{p}{)}
    \PY{n}{cc\PYZus{}list\PYZus{}chunglu}\PY{o}{.}\PY{n}{append}\PY{p}{(}\PY{n}{nx}\PY{o}{.}\PY{n}{average\PYZus{}clustering}\PY{p}{(}\PY{p}{(}\PY{n}{chung\PYZus{}lu\PYZus{}model\PYZus{}from\PYZus{}degree\PYZus{}sequence}\PY{p}{(}\PY{n}{degrees}\PY{p}{)}\PY{p}{)}\PY{p}{)}\PY{p}{)}
    \PY{n}{num\PYZus{}edges\PYZus{}list\PYZus{}config}\PY{o}{.}\PY{n}{append}\PY{p}{(}\PY{p}{(}\PY{n}{configuration\PYZus{}model\PYZus{}from\PYZus{}degree\PYZus{}sequence}\PY{p}{(}\PY{n}{degrees}\PY{p}{)}\PY{p}{)}\PY{o}{.}\PY{n}{number\PYZus{}of\PYZus{}edges}\PY{p}{(}\PY{p}{)}\PY{p}{)}
    \PY{n}{num\PYZus{}edges\PYZus{}list\PYZus{}chunglu}\PY{o}{.}\PY{n}{append}\PY{p}{(}\PY{p}{(}\PY{n}{chung\PYZus{}lu\PYZus{}model\PYZus{}from\PYZus{}degree\PYZus{}sequence}\PY{p}{(}\PY{n}{degrees}\PY{p}{)}\PY{p}{)}\PY{o}{.}\PY{n}{number\PYZus{}of\PYZus{}edges}\PY{p}{(}\PY{p}{)}\PY{p}{)}
    \PY{n}{configs} \PY{o}{+}\PY{o}{=} \PY{p}{(}\PY{n}{nx}\PY{o}{.}\PY{n}{adjacency\PYZus{}matrix}\PY{p}{(}\PY{n}{configuration\PYZus{}model\PYZus{}from\PYZus{}degree\PYZus{}sequence}\PY{p}{(}\PY{n}{degrees}\PY{p}{)}\PY{p}{)}\PY{p}{)}
    \PY{n}{chung\PYZus{}lus} \PY{o}{+}\PY{o}{=} \PY{p}{(}\PY{n}{nx}\PY{o}{.}\PY{n}{adjacency\PYZus{}matrix}\PY{p}{(}\PY{n}{chung\PYZus{}lu\PYZus{}model\PYZus{}from\PYZus{}degree\PYZus{}sequence}\PY{p}{(}\PY{n}{degrees}\PY{p}{)}\PY{p}{)}\PY{p}{)}
    \PY{n}{model}\PY{o}{=}\PY{n}{configuration\PYZus{}model\PYZus{}from\PYZus{}degree\PYZus{}sequence}\PY{p}{(}\PY{n}{degrees}\PY{p}{)}
    \PY{n}{pos} \PY{o}{=} \PY{n}{nx}\PY{o}{.}\PY{n}{spring\PYZus{}layout}\PY{p}{(}\PY{n}{model}\PY{p}{)}

\PY{n}{configs} \PY{o}{=} \PY{n}{configs}\PY{o}{/}\PY{l+m+mi}{1000}
\PY{n}{chung\PYZus{}lus} \PY{o}{=} \PY{n}{chung\PYZus{}lus}\PY{o}{/}\PY{l+m+mi}{1000}

\PY{n+nb}{print}\PY{p}{(}\PY{l+s+s2}{\PYZdq{}}\PY{l+s+s2}{Depicting the configuration nodel graphically to show any self\PYZhy{}loops that can occur: }\PY{l+s+s2}{\PYZdq{}}\PY{p}{)}
\PY{n}{fig2} \PY{o}{=} \PY{n}{plt}\PY{o}{.}\PY{n}{plot}\PY{p}{(}\PY{p}{)}
\PY{n}{nx}\PY{o}{.}\PY{n}{draw\PYZus{}networkx}\PY{p}{(}\PY{n}{model}\PY{p}{,} \PY{n}{pos}\PY{p}{,} \PY{n}{node\PYZus{}size}\PY{o}{=}\PY{l+m+mi}{150}\PY{p}{,} \PY{n}{edgecolors}\PY{o}{=}\PY{l+s+s1}{\PYZsq{}}\PY{l+s+s1}{.2}\PY{l+s+s1}{\PYZsq{}}\PY{p}{,} \PY{n}{with\PYZus{}labels}\PY{o}{=}\PY{k+kc}{True}\PY{p}{,} \PY{n}{font\PYZus{}size}\PY{o}{=}\PY{l+s+s1}{\PYZsq{}}\PY{l+s+s1}{x\PYZhy{}small}\PY{l+s+s1}{\PYZsq{}}\PY{p}{,} \PY{n}{edge\PYZus{}color}\PY{o}{=}\PY{l+s+s1}{\PYZsq{}}\PY{l+s+s1}{.6}\PY{l+s+s1}{\PYZsq{}}\PY{p}{)}


\PY{n}{fig}\PY{p}{,} \PY{n}{ax} \PY{o}{=} \PY{n}{plt}\PY{o}{.}\PY{n}{subplots}\PY{p}{(}\PY{l+m+mi}{1}\PY{p}{,}\PY{l+m+mi}{2}\PY{p}{,}\PY{n}{figsize}\PY{o}{=}\PY{p}{(}\PY{l+m+mi}{8}\PY{p}{,}\PY{l+m+mf}{3.5}\PY{p}{)}\PY{p}{,}\PY{n}{dpi}\PY{o}{=}\PY{l+m+mi}{200}\PY{p}{)}

\PY{n}{im1} \PY{o}{=} \PY{n}{ax}\PY{p}{[}\PY{l+m+mi}{0}\PY{p}{]}\PY{o}{.}\PY{n}{imshow}\PY{p}{(}\PY{n}{configs}\PY{o}{.}\PY{n}{todense}\PY{p}{(}\PY{p}{)}\PY{p}{,}\PY{n}{cmap}\PY{o}{=}\PY{l+s+s1}{\PYZsq{}}\PY{l+s+s1}{viridis}\PY{l+s+s1}{\PYZsq{}}\PY{p}{,}\PY{n}{vmin} \PY{o}{=} \PY{l+m+mi}{0}\PY{p}{,} \PY{n}{vmax} \PY{o}{=} \PY{l+m+mi}{1}\PY{p}{)}
\PY{n}{ax}\PY{p}{[}\PY{l+m+mi}{0}\PY{p}{]}\PY{o}{.}\PY{n}{set\PYZus{}title}\PY{p}{(}\PY{l+s+s2}{\PYZdq{}}\PY{l+s+s2}{Configuration Model Adjacency Matrix Heat Map of 1000 Iterations}\PY{l+s+s2}{\PYZdq{}}\PY{p}{)}
\PY{n}{ax}\PY{p}{[}\PY{l+m+mi}{0}\PY{p}{]}\PY{o}{.}\PY{n}{grid}\PY{p}{(}\PY{k+kc}{True}\PY{p}{)}
\PY{n}{ax}\PY{p}{[}\PY{l+m+mi}{0}\PY{p}{]}\PY{o}{.}\PY{n}{autoscale}\PY{p}{(}\PY{k+kc}{False}\PY{p}{)}

\PY{n}{im2} \PY{o}{=} \PY{n}{ax}\PY{p}{[}\PY{l+m+mi}{1}\PY{p}{]}\PY{o}{.}\PY{n}{imshow}\PY{p}{(}\PY{n}{chung\PYZus{}lus}\PY{o}{.}\PY{n}{todense}\PY{p}{(}\PY{p}{)}\PY{p}{,}\PY{n}{cmap}\PY{o}{=}\PY{l+s+s1}{\PYZsq{}}\PY{l+s+s1}{viridis}\PY{l+s+s1}{\PYZsq{}}\PY{p}{,}\PY{n}{vmin} \PY{o}{=} \PY{l+m+mi}{0}\PY{p}{,} \PY{n}{vmax} \PY{o}{=} \PY{l+m+mi}{1}\PY{p}{)}
\PY{n}{ax}\PY{p}{[}\PY{l+m+mi}{1}\PY{p}{]}\PY{o}{.}\PY{n}{set\PYZus{}title}\PY{p}{(}\PY{l+s+s2}{\PYZdq{}}\PY{l+s+s2}{Chung Lu Model Adjacency Matrix Heat Map of 1000 Iterations}\PY{l+s+s2}{\PYZdq{}}\PY{p}{)}
\PY{n}{ax}\PY{p}{[}\PY{l+m+mi}{1}\PY{p}{]}\PY{o}{.}\PY{n}{grid}\PY{p}{(}\PY{k+kc}{True}\PY{p}{)}
\PY{n}{ax}\PY{p}{[}\PY{l+m+mi}{1}\PY{p}{]}\PY{o}{.}\PY{n}{autoscale}\PY{p}{(}\PY{k+kc}{False}\PY{p}{)}


\PY{n}{fig}\PY{o}{.}\PY{n}{tight\PYZus{}layout}\PY{p}{(}\PY{p}{)}
\PY{n}{fig}\PY{o}{.}\PY{n}{colorbar}\PY{p}{(}\PY{n}{im1}\PY{p}{)}
\PY{n}{fig}\PY{o}{.}\PY{n}{colorbar}\PY{p}{(}\PY{n}{im2}\PY{p}{)}
\PY{n}{plt}\PY{o}{.}\PY{n}{show}\PY{p}{(}\PY{p}{)}

\PY{c+c1}{\PYZsh{} (c) (ii)}
\PY{c+c1}{\PYZsh{}For the metrics, are there any specific ones we want here?}
\PY{c+c1}{\PYZsh{} Could I do just degree distribution? and total of edges?}

\PY{n}{N}\PY{o}{=}\PY{n+nb}{len}\PY{p}{(}\PY{n}{degrees}\PY{p}{)}

\PY{c+c1}{\PYZsh{}\PYZhy{}\PYZhy{}\PYZhy{}Plotting the Total Number of Edges \PYZsq{}L\PYZsq{}}
\PY{c+c1}{\PYZsh{}Starting the plotting process}
\PY{n}{fig}\PY{p}{,} \PY{n}{ax} \PY{o}{=} \PY{n}{plt}\PY{o}{.}\PY{n}{subplots}\PY{p}{(}\PY{p}{)}\PY{c+c1}{\PYZsh{}Note: fig is the plot container, while ax (i.e., the \PYZsq{}axes\PYZsq{}) is the plot itself}
\PY{n}{ax}\PY{o}{.}\PY{n}{hist}\PY{p}{(}\PY{n}{num\PYZus{}edges\PYZus{}list\PYZus{}config}\PY{p}{,}\PY{n}{label}\PY{o}{=}\PY{l+s+s2}{\PYZdq{}}\PY{l+s+s2}{Configuration Model Sample}\PY{l+s+s2}{\PYZdq{}}\PY{p}{,}\PY{n}{density}\PY{o}{=}\PY{k+kc}{True}\PY{p}{)}\PY{c+c1}{\PYZsh{}Plotting the configuration model graph empirical data}
\PY{n}{ax}\PY{o}{.}\PY{n}{hist}\PY{p}{(}\PY{n}{num\PYZus{}edges\PYZus{}list\PYZus{}chunglu}\PY{p}{,}\PY{n}{label}\PY{o}{=}\PY{l+s+s2}{\PYZdq{}}\PY{l+s+s2}{Chung Lu Model Sample}\PY{l+s+s2}{\PYZdq{}}\PY{p}{,}\PY{n}{density}\PY{o}{=}\PY{k+kc}{True}\PY{p}{)}\PY{c+c1}{\PYZsh{}Plotting the Chung Lu model graph empirical data}
\PY{c+c1}{\PYZsh{}ax.plot(L,P\PYZus{}l, color = \PYZsq{}r\PYZsq{},label=\PYZdq{}Theoretical Distribution\PYZdq{}) \PYZsh{}Plotting the Theoretical}
\PY{c+c1}{\PYZsh{}Specifying figure details below:}
\PY{n}{handles}\PY{p}{,} \PY{n}{labels} \PY{o}{=} \PY{n}{ax}\PY{o}{.}\PY{n}{get\PYZus{}legend\PYZus{}handles\PYZus{}labels}\PY{p}{(}\PY{p}{)}
\PY{n}{ax}\PY{o}{.}\PY{n}{legend}\PY{p}{(}\PY{n}{handles}\PY{o}{=}\PY{n}{handles}\PY{p}{,} \PY{n}{labels}\PY{o}{=}\PY{n}{labels}\PY{p}{,}\PY{n}{loc}\PY{o}{=}\PY{l+s+s2}{\PYZdq{}}\PY{l+s+s2}{upper right}\PY{l+s+s2}{\PYZdq{}}\PY{p}{,}\PY{n}{fontsize}\PY{o}{=}\PY{l+s+s2}{\PYZdq{}}\PY{l+s+s2}{small}\PY{l+s+s2}{\PYZdq{}}\PY{p}{)}
\PY{n}{ax}\PY{o}{.}\PY{n}{set\PYZus{}title}\PY{p}{(}\PY{l+s+s2}{\PYZdq{}}\PY{l+s+s2}{Configuration Model vs. Chung Lu Model The Total Number of Edges: }\PY{l+s+s2}{\PYZsq{}}\PY{l+s+s2}{L}\PY{l+s+s2}{\PYZsq{}}\PY{l+s+s2}{\PYZdq{}}\PY{p}{)}
\PY{n}{ax}\PY{o}{.}\PY{n}{set\PYZus{}xlabel}\PY{p}{(}\PY{l+s+s2}{\PYZdq{}}\PY{l+s+s2}{L}\PY{l+s+s2}{\PYZdq{}}\PY{p}{)}
\PY{n}{ax}\PY{o}{.}\PY{n}{set\PYZus{}ylabel}\PY{p}{(}\PY{l+s+s2}{\PYZdq{}}\PY{l+s+s2}{Probability P\PYZus{}L}\PY{l+s+s2}{\PYZdq{}}\PY{p}{)}
\PY{n}{true\PYZus{}min}\PY{o}{=}\PY{n+nb}{min}\PY{p}{(}\PY{n+nb}{min}\PY{p}{(}\PY{n}{num\PYZus{}edges\PYZus{}list\PYZus{}config}\PY{p}{)}\PY{p}{,} \PY{n+nb}{min}\PY{p}{(}\PY{n}{num\PYZus{}edges\PYZus{}list\PYZus{}chunglu}\PY{p}{)}\PY{p}{)}
\PY{n}{true\PYZus{}max}\PY{o}{=}\PY{n+nb}{max}\PY{p}{(}\PY{n+nb}{max}\PY{p}{(}\PY{n}{num\PYZus{}edges\PYZus{}list\PYZus{}config}\PY{p}{)}\PY{p}{,} \PY{n+nb}{max}\PY{p}{(}\PY{n}{num\PYZus{}edges\PYZus{}list\PYZus{}chunglu}\PY{p}{)}\PY{p}{)}
\PY{n}{ax}\PY{o}{.}\PY{n}{set\PYZus{}xlim}\PY{p}{(}\PY{n}{true\PYZus{}min}\PY{p}{,}\PY{n}{true\PYZus{}max}\PY{p}{)}
\PY{n}{plt}\PY{o}{.}\PY{n}{show}\PY{p}{(}\PY{p}{)}

\PY{c+c1}{\PYZsh{}\PYZhy{}\PYZhy{}\PYZhy{}Plotting the Average Clustering Coefficient \PYZsq{}CC\PYZsq{}}
\PY{c+c1}{\PYZsh{}Starting the plotting process}
\PY{n}{fig}\PY{p}{,} \PY{n}{ax} \PY{o}{=} \PY{n}{plt}\PY{o}{.}\PY{n}{subplots}\PY{p}{(}\PY{p}{)}\PY{c+c1}{\PYZsh{}Note: fig is the plot container, while ax (i.e., the \PYZsq{}axes\PYZsq{}) is the plot itself}
\PY{n}{ax}\PY{o}{.}\PY{n}{hist}\PY{p}{(}\PY{n}{cc\PYZus{}list\PYZus{}config}\PY{p}{,}\PY{n}{label}\PY{o}{=}\PY{l+s+s2}{\PYZdq{}}\PY{l+s+s2}{Configuration Model Sample}\PY{l+s+s2}{\PYZdq{}}\PY{p}{,}\PY{n}{density}\PY{o}{=}\PY{k+kc}{True}\PY{p}{)}\PY{c+c1}{\PYZsh{}Plotting the configuration model graph empirical data}
\PY{n}{ax}\PY{o}{.}\PY{n}{hist}\PY{p}{(}\PY{n}{cc\PYZus{}list\PYZus{}chunglu}\PY{p}{,}\PY{n}{label}\PY{o}{=}\PY{l+s+s2}{\PYZdq{}}\PY{l+s+s2}{Chung Lu Model Sample}\PY{l+s+s2}{\PYZdq{}}\PY{p}{,}\PY{n}{density}\PY{o}{=}\PY{k+kc}{True}\PY{p}{)}\PY{c+c1}{\PYZsh{}Plotting the Chung Lu model graph empirical data}
\PY{c+c1}{\PYZsh{}ax.plot(L,P\PYZus{}l, color = \PYZsq{}r\PYZsq{},label=\PYZdq{}Theoretical Distribution\PYZdq{}) \PYZsh{}Plotting the Theoretical}
\PY{c+c1}{\PYZsh{}Specifying figure details below:}
\PY{n}{handles}\PY{p}{,} \PY{n}{labels} \PY{o}{=} \PY{n}{ax}\PY{o}{.}\PY{n}{get\PYZus{}legend\PYZus{}handles\PYZus{}labels}\PY{p}{(}\PY{p}{)}
\PY{n}{ax}\PY{o}{.}\PY{n}{legend}\PY{p}{(}\PY{n}{handles}\PY{o}{=}\PY{n}{handles}\PY{p}{,} \PY{n}{labels}\PY{o}{=}\PY{n}{labels}\PY{p}{,}\PY{n}{loc}\PY{o}{=}\PY{l+s+s2}{\PYZdq{}}\PY{l+s+s2}{upper right}\PY{l+s+s2}{\PYZdq{}}\PY{p}{,}\PY{n}{fontsize}\PY{o}{=}\PY{l+s+s2}{\PYZdq{}}\PY{l+s+s2}{small}\PY{l+s+s2}{\PYZdq{}}\PY{p}{)}
\PY{n}{ax}\PY{o}{.}\PY{n}{set\PYZus{}title}\PY{p}{(}\PY{l+s+s2}{\PYZdq{}}\PY{l+s+s2}{Configuration Model vs. Chung Lu Model Average Clustering Coefficient }\PY{l+s+s2}{\PYZsq{}}\PY{l+s+s2}{CC}\PY{l+s+s2}{\PYZsq{}}\PY{l+s+s2}{\PYZdq{}}\PY{p}{)}
\PY{n}{ax}\PY{o}{.}\PY{n}{set\PYZus{}xlabel}\PY{p}{(}\PY{l+s+s2}{\PYZdq{}}\PY{l+s+s2}{Average Clustering Coefficient}\PY{l+s+s2}{\PYZdq{}}\PY{p}{)}
\PY{n}{ax}\PY{o}{.}\PY{n}{set\PYZus{}ylabel}\PY{p}{(}\PY{l+s+s2}{\PYZdq{}}\PY{l+s+s2}{Probability P\PYZus{}CC}\PY{l+s+s2}{\PYZdq{}}\PY{p}{)}
\PY{n}{true\PYZus{}min}\PY{o}{=}\PY{n+nb}{min}\PY{p}{(}\PY{n+nb}{min}\PY{p}{(}\PY{n}{cc\PYZus{}list\PYZus{}config}\PY{p}{)}\PY{p}{,} \PY{n+nb}{min}\PY{p}{(}\PY{n}{cc\PYZus{}list\PYZus{}chunglu}\PY{p}{)}\PY{p}{)}
\PY{n}{true\PYZus{}max}\PY{o}{=}\PY{n+nb}{max}\PY{p}{(}\PY{n+nb}{max}\PY{p}{(}\PY{n}{cc\PYZus{}list\PYZus{}config}\PY{p}{)}\PY{p}{,} \PY{n+nb}{max}\PY{p}{(}\PY{n}{cc\PYZus{}list\PYZus{}chunglu}\PY{p}{)}\PY{p}{)}
\PY{n}{ax}\PY{o}{.}\PY{n}{set\PYZus{}xlim}\PY{p}{(}\PY{n}{true\PYZus{}min}\PY{p}{,}\PY{n}{true\PYZus{}max}\PY{p}{)}
\PY{n}{plt}\PY{o}{.}\PY{n}{show}\PY{p}{(}\PY{p}{)}
\end{Verbatim}
\end{tcolorbox}

    \begin{Verbatim}[commandchars=\\\{\}]
[(0, 5), (1, 3), (2, 3), (3, 4), (4, 1), (5, 2), (6, 2), (7, 4), (8, 2), (9, 0),
(10, 0), (11, 0), (12, 0), (13, 0), (14, 0)]
Depicting the configuration nodel graphically to show any self-loops that can
occur:
    \end{Verbatim}

    \begin{center}
    \adjustimage{max size={0.9\linewidth}{0.9\paperheight}}{assignment02_files/assignment02_1_1.png}
    \end{center}
    { \hspace*{\fill} \\}
    
    \begin{center}
    \adjustimage{max size={0.9\linewidth}{0.9\paperheight}}{assignment02_files/assignment02_1_2.png}
    \end{center}
    { \hspace*{\fill} \\}
    
    \begin{center}
    \adjustimage{max size={0.9\linewidth}{0.9\paperheight}}{assignment02_files/assignment02_1_3.png}
    \end{center}
    { \hspace*{\fill} \\}
    
    \begin{center}
    \adjustimage{max size={0.9\linewidth}{0.9\paperheight}}{assignment02_files/assignment02_1_4.png}
    \end{center}
    { \hspace*{\fill} \\}
    
    Discussing the Results: My heat map comparing the two models shows the
general adjacency matrix value for each pair of nodes. Notably we see
the Chung Lu model guaranteed no self-loops, while the Configuration
Model had some. The general trend of the higher degree nodes (i.e., in
this case node 0,1,2, etc., as I enforced and provided a sorted list from
highest degree to lowest as input) remained the region with the highest
overall adjacency matrix average values amongst the 1000 iterations.
This indicates they had a high probability to be connected amongst one
another, and in turn were thus often connected to one another. This is
expected as the high degree nodes are the most likely for any random
node to be connected to. In the case of Chung Lu however some of these
pairs appeared to be connected close to if not 100\% of the time, in
contrast to the Configuration where these high-degree to high-degree
pairs (e.g., node 0 to node 1), were the highest but did not occur 100\%
of the time across iterations.

As for my two metrics I chose for part c, I did the total number of
edges `L', and the average clustering coefficient. Here we can see the
total number of edges, is approximately the same across graphs generated
by both models, with both models coalescing around the same average
total number of edges, of roughly \textasciitilde17 or so, albeit at
different probability levels. The average clustering coefficient on the
other hand is much closer to 0 in the case of the configuration model
while the Chung Lu model is relatively higher than 0. This is most
likely due to the presence of self-loops and multi-edges that can occur
in the configuration model. This artifact of the self-loops and
multi-edges results in less clustering amongst nodes, relative to the
Chung Lu model's differring approach.


\begin{problemstmt}
\textbf{Problem 2 (20 points)}

In this question we will gain experience with community detection.

\begin{enumerate}
\def\labelenumi{(\alph{enumi})}
\tightlist
\item
  Finding communities (5 points).
\end{enumerate}

Download a network dataset from https://networkrepository.com/ with
between roughly 1,000 \textless{} N \textless{} 5,000 nodes (select a
different dataset from the one you used in assignment 1). Select a
modularity maximization algorithm to detect communities in this graph.

\begin{enumerate}
\def\labelenumi{(\roman{enumi})}
\item
  Visualize your network, coloring the nodes by community membership.
\item
  Report the modularity of the partition your algorithm found.
\item
  Report the number of communities detected.
\item
  Store your partition as a dictionary in the form of \{node id:
  community id\}.
\end{enumerate}

\begin{enumerate}
\def\labelenumi{(\alph{enumi})}
\setcounter{enumi}{1}
\tightlist
\item
  Null models (5 points).
\end{enumerate}

Create a degree-preserving randomization of your network, and run the
community detection algorithm you chose in part a on your randomized
graph. Store each partition as a dictionary in the form of \{node id:
community id\}.

\begin{enumerate}
\def\labelenumi{(\alph{enumi})}
\setcounter{enumi}{2}
\tightlist
\item
  Plotting a null model comparison (5 points).
\end{enumerate}

You should now have two dictionaries, corresponding to the partitions
from parts a and b. Create a figure with two subplots: In each, plot the
community size distribution of each partition. Your figure should be
publication-quality, with well-labeled, consistent axes, log-scaled
wherever relevant, etc.

\begin{enumerate}
\def\labelenumi{(\alph{enumi})}
\setcounter{enumi}{3}
\tightlist
\item
  Analysis (5 points).
\end{enumerate}

For the different partitions computed above, discuss how you would
approach comparing this functions quantitatively. For example, imagine a
function f(p1,p2) that takes partitions p1 and p2 as input and returns a
number corresponding to how similar or dissimilar the partitions are.
What properties should this function f have? For this question, base
your answer on a technique used in the literature, after a brief search
about methods for quantifying differences between graph partitions. You
do not need to implement the function, a high-level discussion is
sufficient.
\end{problemstmt}

    \begin{tcolorbox}[breakable, size=fbox, boxrule=1pt, pad at break*=1mm,colback=cellbackground, colframe=cellborder]
\prompt{In}{incolor}{18}{\boxspacing}
\begin{Verbatim}[commandchars=\\\{\}]
\PY{k+kn}{from}\PY{+w}{ }\PY{n+nn}{scipy}\PY{n+nn}{.}\PY{n+nn}{io}\PY{+w}{ }\PY{k+kn}{import} \PY{n}{mmread}

\PY{n}{a} \PY{o}{=} \PY{n}{mmread}\PY{p}{(}\PY{l+s+s1}{\PYZsq{}}\PY{l+s+s1}{ca\PYZhy{}CSphd.mtx}\PY{l+s+s1}{\PYZsq{}}\PY{p}{)}
\PY{n}{phd\PYZus{}collabs} \PY{o}{=}\PY{n}{nx}\PY{o}{.}\PY{n}{DiGraph}\PY{p}{(}\PY{n}{a}\PY{p}{)}\PY{o}{.}\PY{n}{to\PYZus{}undirected}\PY{p}{(}\PY{p}{)}
\end{Verbatim}
\end{tcolorbox}

    \begin{Verbatim}[commandchars=\\\{\}]
/tmp/ipykernel\_156064/3807146788.py:3: DeprecationWarning: The default value for
`spmatrix` is changing to `False` in v1.20.
             That means the default return type will be a sparse array.
             Unless you use * instead of @, ** for matrix power, or you depend
             on 2D shapes from e.g. `A.sum(axis=0)` it may not matter to you.
             See the spmatrix to sparray migration guide for details.
https://docs.scipy.org/doc/scipy/reference/sparse.migration\_to\_sparray.html

  a = mmread('ca-CSphd.mtx')
    \end{Verbatim}

    \begin{tcolorbox}[breakable, size=fbox, boxrule=1pt, pad at break*=1mm,colback=cellbackground, colframe=cellborder]
\prompt{In}{incolor}{11}{\boxspacing}
\end{tcolorbox}

    \begin{tcolorbox}[breakable, size=fbox, boxrule=1pt, pad at break*=1mm,colback=cellbackground, colframe=cellborder]
\prompt{In}{incolor}{ }{\boxspacing}
\begin{Verbatim}[commandchars=\\\{\}]
\PY{k+kn}{import}\PY{+w}{ }\PY{n+nn}{networkx}\PY{+w}{ }\PY{k}{as}\PY{+w}{ }\PY{n+nn}{nx}
\PY{k+kn}{import}\PY{+w}{ }\PY{n+nn}{numpy}\PY{+w}{ }\PY{k}{as}\PY{+w}{ }\PY{n+nn}{np}
\PY{k+kn}{from}\PY{+w}{ }\PY{n+nn}{collections}\PY{+w}{ }\PY{k+kn}{import} \PY{n}{defaultdict}
\PY{k+kn}{from}\PY{+w}{ }\PY{n+nn}{random}\PY{+w}{ }\PY{k+kn}{import} \PY{n}{shuffle}
\PY{k+kn}{import}\PY{+w}{ }\PY{n+nn}{matplotlib}\PY{n+nn}{.}\PY{n+nn}{pyplot}\PY{+w}{ }\PY{k}{as}\PY{+w}{ }\PY{n+nn}{plt}
\PY{k+kn}{import}\PY{+w}{ }\PY{n+nn}{community}

\PY{n}{pos} \PY{o}{=} \PY{n}{nx}\PY{o}{.}\PY{n}{fruchterman\PYZus{}reingold\PYZus{}layout}\PY{p}{(}\PY{n}{phd\PYZus{}collabs}\PY{p}{)}
\PY{n}{partition} \PY{o}{=} \PY{n}{community}\PY{o}{.}\PY{n}{best\PYZus{}partition}\PY{p}{(}\PY{n}{phd\PYZus{}collabs}\PY{p}{)}

\PY{c+c1}{\PYZsh{} Output the best modularity score}
\PY{n+nb}{print}\PY{p}{(}\PY{l+s+s2}{\PYZdq{}}\PY{l+s+s2}{Best Modularity: }\PY{l+s+s2}{\PYZdq{}}\PY{p}{,} \PY{n}{community}\PY{o}{.}\PY{n}{modularity}\PY{p}{(}\PY{n}{partition}\PY{p}{,} \PY{n}{phd\PYZus{}collabs}\PY{p}{)}\PY{p}{)}

\PY{c+c1}{\PYZsh{} Output the Number of Communities Detected \PYZob{}node\PYZus{}id: community\PYZus{}id\PYZcb{}}
\PY{n+nb}{print}\PY{p}{(}\PY{l+s+s2}{\PYZdq{}}\PY{l+s+s2}{Number of Communities Detected: }\PY{l+s+s2}{\PYZdq{}}\PY{p}{,}\PY{n+nb}{len}\PY{p}{(}\PY{n}{np}\PY{o}{.}\PY{n}{unique}\PY{p}{(}\PY{n+nb}{list}\PY{p}{(}\PY{n}{partition}\PY{o}{.}\PY{n}{values}\PY{p}{(}\PY{p}{)}\PY{p}{)}\PY{p}{)}\PY{p}{)}\PY{p}{)}

\PY{c+c1}{\PYZsh{} Storing and showing partition as dictionary}
\PY{n}{phd\PYZus{}collabs\PYZus{}parts} \PY{o}{=} \PY{n+nb}{dict}\PY{p}{(}\PY{n}{partition}\PY{p}{)}
\PY{n+nb}{print}\PY{p}{(}\PY{l+s+s2}{\PYZdq{}}\PY{l+s+s2}{Partition Dictionary: }\PY{l+s+s2}{\PYZdq{}}\PY{p}{,} \PY{n}{phd\PYZus{}collabs\PYZus{}parts}\PY{p}{)}
\PY{n}{colors\PYZus{}sns} \PY{o}{=} \PY{n}{plt}\PY{o}{.}\PY{n}{cm}\PY{o}{.}\PY{n}{tab20b}\PY{p}{(}\PY{n}{np}\PY{o}{.}\PY{n}{linspace}\PY{p}{(}\PY{l+m+mi}{0}\PY{p}{,}\PY{l+m+mi}{1}\PY{p}{,}\PY{n+nb}{len}\PY{p}{(}\PY{n}{np}\PY{o}{.}\PY{n}{unique}\PY{p}{(}\PY{n+nb}{list}\PY{p}{(}\PY{n}{partition}\PY{o}{.}\PY{n}{values}\PY{p}{(}\PY{p}{)}\PY{p}{)}\PY{p}{)}\PY{p}{)}\PY{p}{)}\PY{p}{)}
\PY{n}{colors\PYZus{}nodes} \PY{o}{=} \PY{p}{[}\PY{n}{colors\PYZus{}sns}\PY{p}{[}\PY{n}{partition}\PY{p}{[}\PY{n}{node}\PY{p}{]}\PY{p}{]} \PY{k}{for} \PY{n}{node} \PY{o+ow}{in} \PY{n}{phd\PYZus{}collabs}\PY{o}{.}\PY{n}{nodes}\PY{p}{(}\PY{p}{)}\PY{p}{]}

\PY{n}{fig}\PY{p}{,} \PY{n}{ax} \PY{o}{=} \PY{n}{plt}\PY{o}{.}\PY{n}{subplots}\PY{p}{(}\PY{l+m+mi}{1}\PY{p}{,}\PY{l+m+mi}{1}\PY{p}{,}\PY{n}{figsize}\PY{o}{=}\PY{p}{(}\PY{l+m+mi}{7}\PY{p}{,}\PY{l+m+mi}{5}\PY{p}{)}\PY{p}{,}\PY{n}{dpi}\PY{o}{=}\PY{l+m+mi}{100}\PY{p}{)}

\PY{n}{nx}\PY{o}{.}\PY{n}{draw}\PY{p}{(}\PY{n}{phd\PYZus{}collabs}\PY{p}{,} \PY{n}{pos}\PY{o}{=}\PY{n}{pos}\PY{p}{,} \PY{n}{node\PYZus{}size}\PY{o}{=}\PY{l+m+mi}{150}\PY{p}{,} \PY{n}{edgecolors}\PY{o}{=}\PY{l+s+s1}{\PYZsq{}}\PY{l+s+s1}{.2}\PY{l+s+s1}{\PYZsq{}}\PY{p}{,} \PY{n}{node\PYZus{}color}\PY{o}{=}\PY{n}{colors\PYZus{}nodes}\PY{p}{,} \PY{n}{ax}\PY{o}{=}\PY{n}{ax}\PY{p}{)}

\PY{n}{plt}\PY{o}{.}\PY{n}{show}\PY{p}{(}\PY{p}{)}
\end{Verbatim}
\end{tcolorbox}

    \begin{Verbatim}[commandchars=\\\{\}]
Best Modularity:  0.9608196261068275 0.9608196261064889
Number of Communities Detected:  189
Partition Dictionary:  \{0: 19, 1: 19, 2: 1, 3: 1, 4: 2, 5: 2, 6: 3, 7: 3, 8: 4,
\}
    \end{Verbatim}

    \begin{center}
    \adjustimage{max size={0.9\linewidth}{0.9\paperheight}}{assignment02_files/assignment02_6_1.png}
    \end{center}
    { \hspace*{\fill} \\}
    
    \begin{tcolorbox}[breakable, size=fbox, boxrule=1pt, pad at break*=1mm,colback=cellbackground, colframe=cellborder]
\prompt{In}{incolor}{41}{\boxspacing}
\begin{Verbatim}[commandchars=\\\{\}]
\PY{c+c1}{\PYZsh{}(b)}
\PY{n}{degree\PYZus{}dict} \PY{o}{=} \PY{n+nb}{dict}\PY{p}{(}\PY{n}{phd\PYZus{}collabs}\PY{o}{.}\PY{n}{degree}\PY{p}{(}\PY{p}{)}\PY{p}{)}
\PY{n}{degree\PYZus{}sequence} \PY{o}{=} \PY{n+nb}{list}\PY{p}{(}\PY{n}{degree\PYZus{}dict}\PY{o}{.}\PY{n}{values}\PY{p}{(}\PY{p}{)}\PY{p}{)}
\PY{n}{null\PYZus{}phd\PYZus{}collabs}\PY{o}{=}\PY{n}{chung\PYZus{}lu\PYZus{}model\PYZus{}from\PYZus{}degree\PYZus{}sequence}\PY{p}{(}\PY{n}{degree\PYZus{}sequence}\PY{p}{)}
\PY{n}{pos} \PY{o}{=} \PY{n}{nx}\PY{o}{.}\PY{n}{fruchterman\PYZus{}reingold\PYZus{}layout}\PY{p}{(}\PY{n}{null\PYZus{}phd\PYZus{}collabs}\PY{p}{)}
\PY{n}{partition} \PY{o}{=} \PY{n}{community}\PY{o}{.}\PY{n}{best\PYZus{}partition}\PY{p}{(}\PY{n}{null\PYZus{}phd\PYZus{}collabs}\PY{p}{)}
\PY{c+c1}{\PYZsh{} Output the best modularity score}
\PY{n+nb}{print}\PY{p}{(}\PY{l+s+s2}{\PYZdq{}}\PY{l+s+s2}{Best Modularity: }\PY{l+s+s2}{\PYZdq{}}\PY{p}{,}\PY{n}{community}\PY{o}{.}\PY{n}{modularity}\PY{p}{(}\PY{n}{partition}\PY{p}{,} \PY{n}{null\PYZus{}phd\PYZus{}collabs}\PY{p}{)}\PY{p}{)}

\PY{c+c1}{\PYZsh{} Output the Number of Communities Detected \PYZob{}node\PYZus{}id: community\PYZus{}id\PYZcb{}}
\PY{n+nb}{print}\PY{p}{(}\PY{l+s+s2}{\PYZdq{}}\PY{l+s+s2}{Number of Communities Detected: }\PY{l+s+s2}{\PYZdq{}}\PY{p}{,}\PY{n+nb}{len}\PY{p}{(}\PY{n}{np}\PY{o}{.}\PY{n}{unique}\PY{p}{(}\PY{n+nb}{list}\PY{p}{(}\PY{n}{partition}\PY{o}{.}\PY{n}{values}\PY{p}{(}\PY{p}{)}\PY{p}{)}\PY{p}{)}\PY{p}{)}\PY{p}{)}

\PY{c+c1}{\PYZsh{} Storing and showing partition as dictionary}
\PY{n}{null\PYZus{}phd\PYZus{}collabs\PYZus{}parts} \PY{o}{=} \PY{n+nb}{dict}\PY{p}{(}\PY{n}{partition}\PY{p}{)}
\PY{n+nb}{print}\PY{p}{(}\PY{l+s+s2}{\PYZdq{}}\PY{l+s+s2}{Partition Dictionary: }\PY{l+s+s2}{\PYZdq{}}\PY{p}{,} \PY{n}{null\PYZus{}phd\PYZus{}collabs\PYZus{}parts}\PY{p}{)}
\PY{n}{colors\PYZus{}sns} \PY{o}{=} \PY{n}{plt}\PY{o}{.}\PY{n}{cm}\PY{o}{.}\PY{n}{tab20b}\PY{p}{(}\PY{n}{np}\PY{o}{.}\PY{n}{linspace}\PY{p}{(}\PY{l+m+mi}{0}\PY{p}{,}\PY{l+m+mi}{1}\PY{p}{,}\PY{n+nb}{len}\PY{p}{(}\PY{n}{np}\PY{o}{.}\PY{n}{unique}\PY{p}{(}\PY{n+nb}{list}\PY{p}{(}\PY{n}{partition}\PY{o}{.}\PY{n}{values}\PY{p}{(}\PY{p}{)}\PY{p}{)}\PY{p}{)}\PY{p}{)}\PY{p}{)}\PY{p}{)}
\PY{n}{colors\PYZus{}nodes} \PY{o}{=} \PY{p}{[}\PY{n}{colors\PYZus{}sns}\PY{p}{[}\PY{n}{partition}\PY{p}{[}\PY{n}{node}\PY{p}{]}\PY{p}{]} \PY{k}{for} \PY{n}{node} \PY{o+ow}{in} \PY{n}{null\PYZus{}phd\PYZus{}collabs}\PY{o}{.}\PY{n}{nodes}\PY{p}{(}\PY{p}{)}\PY{p}{]}

\PY{n}{fig}\PY{p}{,} \PY{n}{ax} \PY{o}{=} \PY{n}{plt}\PY{o}{.}\PY{n}{subplots}\PY{p}{(}\PY{l+m+mi}{1}\PY{p}{,}\PY{l+m+mi}{1}\PY{p}{,}\PY{n}{figsize}\PY{o}{=}\PY{p}{(}\PY{l+m+mi}{7}\PY{p}{,}\PY{l+m+mi}{5}\PY{p}{)}\PY{p}{,}\PY{n}{dpi}\PY{o}{=}\PY{l+m+mi}{100}\PY{p}{)}

\PY{n}{nx}\PY{o}{.}\PY{n}{draw}\PY{p}{(}\PY{n}{null\PYZus{}phd\PYZus{}collabs}\PY{p}{,} \PY{n}{pos}\PY{o}{=}\PY{n}{pos}\PY{p}{,} \PY{n}{node\PYZus{}size}\PY{o}{=}\PY{l+m+mi}{150}\PY{p}{,} \PY{n}{edgecolors}\PY{o}{=}\PY{l+s+s1}{\PYZsq{}}\PY{l+s+s1}{.2}\PY{l+s+s1}{\PYZsq{}}\PY{p}{,} \PY{n}{node\PYZus{}color}\PY{o}{=}\PY{n}{colors\PYZus{}nodes}\PY{p}{,} \PY{n}{ax}\PY{o}{=}\PY{n}{ax}\PY{p}{)}

\PY{n}{plt}\PY{o}{.}\PY{n}{show}\PY{p}{(}\PY{p}{)}
\end{Verbatim}
\end{tcolorbox}

    \begin{Verbatim}[commandchars=\\\{\}]
Best Modularity:  0.7232526057391191
Number of Communities Detected:  618
Partition Dictionary:  \{0: 0, 1: 38, 2: 17, 3: 3, 4: 30, 5: 26, 6: 17, 7: 7, 8:
\}
    \end{Verbatim}

    \begin{center}
    \adjustimage{max size={0.9\linewidth}{0.9\paperheight}}{assignment02_files/assignment02_7_1.png}
    \end{center}
    { \hspace*{\fill} \\}
    
    \begin{tcolorbox}[breakable, size=fbox, boxrule=1pt, pad at break*=1mm,colback=cellbackground, colframe=cellborder]
\prompt{In}{incolor}{51}{\boxspacing}
\begin{Verbatim}[commandchars=\\\{\}]
\PY{k+kn}{from}\PY{+w}{ }\PY{n+nn}{collections}\PY{+w}{ }\PY{k+kn}{import} \PY{n}{Counter}

\PY{n}{phd\PYZus{}collabs\PYZus{}community\PYZus{}sizes} \PY{o}{=} \PY{n}{Counter}\PY{p}{(}\PY{n}{phd\PYZus{}collabs\PYZus{}parts}\PY{o}{.}\PY{n}{values}\PY{p}{(}\PY{p}{)}\PY{p}{)}
\PY{n}{null\PYZus{}phd\PYZus{}collabs\PYZus{}community\PYZus{}sizes} \PY{o}{=} \PY{n}{Counter}\PY{p}{(}\PY{n}{null\PYZus{}phd\PYZus{}collabs\PYZus{}parts}\PY{o}{.}\PY{n}{values}\PY{p}{(}\PY{p}{)}\PY{p}{)}

\PY{n}{fig}\PY{p}{,} \PY{n}{ax} \PY{o}{=} \PY{n}{plt}\PY{o}{.}\PY{n}{subplots}\PY{p}{(}\PY{l+m+mi}{1}\PY{p}{,}\PY{l+m+mi}{2}\PY{p}{,}\PY{n}{figsize}\PY{o}{=}\PY{p}{(}\PY{l+m+mi}{8}\PY{p}{,}\PY{l+m+mf}{3.5}\PY{p}{)}\PY{p}{,}\PY{n}{dpi}\PY{o}{=}\PY{l+m+mi}{300}\PY{p}{,}\PY{n}{sharey}\PY{o}{=}\PY{k+kc}{True}\PY{p}{,}\PY{n}{sharex}\PY{o}{=}\PY{k+kc}{True}\PY{p}{)}

\PY{c+c1}{\PYZsh{} Plotting the community size distribution of the original dataset}
\PY{n}{org\PYZus{}bins} \PY{o}{=} \PY{n}{np}\PY{o}{.}\PY{n}{geomspace}\PY{p}{(}\PY{l+m+mi}{1}\PY{p}{,}\PY{n}{np}\PY{o}{.}\PY{n}{max}\PY{p}{(}\PY{n+nb}{list}\PY{p}{(}\PY{n}{phd\PYZus{}collabs\PYZus{}community\PYZus{}sizes}\PY{p}{)}\PY{p}{)}\PY{p}{)}\PY{c+c1}{\PYZsh{}Defining the smallest community as an isolated node of 1}
\PY{n}{ax}\PY{p}{[}\PY{l+m+mi}{0}\PY{p}{]}\PY{o}{.}\PY{n}{hist}\PY{p}{(}\PY{n+nb}{list}\PY{p}{(}\PY{n}{phd\PYZus{}collabs\PYZus{}community\PYZus{}sizes}\PY{p}{)}\PY{p}{,}\PY{n}{bins}\PY{o}{=}\PY{n}{org\PYZus{}bins}\PY{p}{,} \PY{n}{density}\PY{o}{=}\PY{k+kc}{True}\PY{p}{,}\PY{n}{color} \PY{o}{=} \PY{l+s+s1}{\PYZsq{}}\PY{l+s+s1}{b}\PY{l+s+s1}{\PYZsq{}}\PY{p}{,}\PY{n}{histtype}\PY{o}{=}\PY{l+s+s1}{\PYZsq{}}\PY{l+s+s1}{step}\PY{l+s+s1}{\PYZsq{}}\PY{p}{)}
\PY{n}{ax}\PY{p}{[}\PY{l+m+mi}{0}\PY{p}{]}\PY{o}{.}\PY{n}{loglog}\PY{p}{(}\PY{p}{)}
\PY{n}{ax}\PY{p}{[}\PY{l+m+mi}{0}\PY{p}{]}\PY{o}{.}\PY{n}{set\PYZus{}title}\PY{p}{(}\PY{l+s+s2}{\PYZdq{}}\PY{l+s+s2}{Original PhD Collaborations Data}\PY{l+s+s2}{\PYZdq{}}\PY{p}{)}
\PY{n}{ax}\PY{p}{[}\PY{l+m+mi}{0}\PY{p}{]}\PY{o}{.}\PY{n}{set\PYZus{}xlabel}\PY{p}{(}\PY{l+s+s2}{\PYZdq{}}\PY{l+s+s2}{Community Size}\PY{l+s+s2}{\PYZdq{}}\PY{p}{)}
\PY{n}{ax}\PY{p}{[}\PY{l+m+mi}{0}\PY{p}{]}\PY{o}{.}\PY{n}{set\PYZus{}ylabel}\PY{p}{(}\PY{l+s+s2}{\PYZdq{}}\PY{l+s+s2}{Probability}\PY{l+s+s2}{\PYZdq{}}\PY{p}{)}
\PY{n}{ax}\PY{p}{[}\PY{l+m+mi}{0}\PY{p}{]}\PY{o}{.}\PY{n}{grid}\PY{p}{(}\PY{k+kc}{True}\PY{p}{)}
\PY{c+c1}{\PYZsh{} Plotting the community size distribution of the Chung Lu Null Model of the dataset}
\PY{n}{null\PYZus{}bins} \PY{o}{=} \PY{n}{np}\PY{o}{.}\PY{n}{geomspace}\PY{p}{(}\PY{l+m+mi}{1}\PY{p}{,}\PY{n}{np}\PY{o}{.}\PY{n}{max}\PY{p}{(}\PY{n+nb}{list}\PY{p}{(}\PY{n}{null\PYZus{}phd\PYZus{}collabs\PYZus{}community\PYZus{}sizes}\PY{p}{)}\PY{p}{)}\PY{p}{)}
\PY{n}{ax}\PY{p}{[}\PY{l+m+mi}{1}\PY{p}{]}\PY{o}{.}\PY{n}{hist}\PY{p}{(}\PY{n+nb}{list}\PY{p}{(}\PY{n}{null\PYZus{}phd\PYZus{}collabs\PYZus{}community\PYZus{}sizes}\PY{p}{)}\PY{p}{,}\PY{n}{bins}\PY{o}{=}\PY{n}{null\PYZus{}bins}\PY{p}{,}\PY{n}{density}\PY{o}{=}\PY{k+kc}{True}\PY{p}{,}\PY{n}{color} \PY{o}{=} \PY{l+s+s1}{\PYZsq{}}\PY{l+s+s1}{b}\PY{l+s+s1}{\PYZsq{}}\PY{p}{,}\PY{n}{histtype}\PY{o}{=}\PY{l+s+s1}{\PYZsq{}}\PY{l+s+s1}{step}\PY{l+s+s1}{\PYZsq{}}\PY{p}{)}
\PY{n}{ax}\PY{p}{[}\PY{l+m+mi}{1}\PY{p}{]}\PY{o}{.}\PY{n}{loglog}\PY{p}{(}\PY{p}{)}
\PY{n}{ax}\PY{p}{[}\PY{l+m+mi}{1}\PY{p}{]}\PY{o}{.}\PY{n}{set\PYZus{}title}\PY{p}{(}\PY{l+s+s2}{\PYZdq{}}\PY{l+s+s2}{Chung Lu Null Model of Original Data}\PY{l+s+s2}{\PYZdq{}}\PY{p}{)}
\PY{n}{ax}\PY{p}{[}\PY{l+m+mi}{1}\PY{p}{]}\PY{o}{.}\PY{n}{set\PYZus{}xlabel}\PY{p}{(}\PY{l+s+s2}{\PYZdq{}}\PY{l+s+s2}{Community Size}\PY{l+s+s2}{\PYZdq{}}\PY{p}{)}
\PY{n}{ax}\PY{p}{[}\PY{l+m+mi}{1}\PY{p}{]}\PY{o}{.}\PY{n}{set\PYZus{}ylabel}\PY{p}{(}\PY{l+s+s2}{\PYZdq{}}\PY{l+s+s2}{Probability}\PY{l+s+s2}{\PYZdq{}}\PY{p}{)}
\PY{n}{ax}\PY{p}{[}\PY{l+m+mi}{1}\PY{p}{]}\PY{o}{.}\PY{n}{tick\PYZus{}params}\PY{p}{(}\PY{n}{labelleft}\PY{o}{=}\PY{k+kc}{True}\PY{p}{)}
\PY{n}{ax}\PY{p}{[}\PY{l+m+mi}{1}\PY{p}{]}\PY{o}{.}\PY{n}{grid}\PY{p}{(}\PY{k+kc}{True}\PY{p}{)}

\PY{n}{plt}\PY{o}{.}\PY{n}{tight\PYZus{}layout}\PY{p}{(}\PY{p}{)}
\end{Verbatim}
\end{tcolorbox}

    \begin{center}
    \adjustimage{max size={0.9\linewidth}{0.9\paperheight}}{assignment02_files/assignment02_8_0.png}
    \end{center}
    { \hspace*{\fill} \\}
    
    Analysis:

In the literature there are a vast number of different approaches to
comparing two partitions: (Fortunato, Santo. ``Community Detection in
Graphs.'' Physics Reports 486, no. 3 (2010): 75--174.
https://doi.org/10.1016/j.physrep.2009.11.002.) provides a really nice
overview of many!

I would likely approach the problem through the classification error
H(X,Y) measure (M. Meilă, D. Heckerman, An experimental comparison of
model-based clustering methods, Mach. Learn. 42 (1) (2001) 9--29.),
which returns back the maximum overlap between the two partitions, in
terms of community assignments of each node. This method seemed the most
intuitive to me. Since, my intial analysis simply assigns a single
community partition per node, this would allow me to compare directly to
see how many nodes were assigned to the same community across the two
partitions. This is good because my dataset reflects PhD Collaborations,
and so being able to distinguish between a person who may be part of 1
community as opposed to definitely apart of 1 community is incredibly
valuable in better understanding my dataset.

The only caveat I think would be that in real-life people tend to have
multiple different communities, each of which they could be involved
with at varying levels, and so by only evaluating the match across a
single community assignment I miss some of that information, however to
analyze assignments to multiple communities requires leveraging a
different community assignment algorithm to the Louvain algorithm I
utilized in my analyses above.

Overall, I think an ideal function, f, to approach this problem would
likely extend the idea of classification error but to a list of all
communities a given node was assigned to across the 2 different
partitions, each with an associated probability score. This would
provide a more robust measure of community assignments per node that
would allow one to see the following properties: whether a node is
generally a part of some subset of communities, or not, to what degree
if so, and how confident each partition was in assigning said
communities. I think these are the primary properties the function f
should have in order to determine well-enough whether a given pair of
partitions are similar or not.

\begin{problemstmt}
    \textbf{Problem 3 (20 points)}

In this question we will revisit degree sequences with a focus on
preferential attachment. We have worked with G(n,p) random graphs
extensively. We have also discussed how most graphs are not random, but
follow some kind of preference for homophily or highly ``desirable''
nodes. One common preference is to connect to high degree nodes,
reinforcing a rich-get-richer dynamic. We will create a network growth
model, where we add new vertices to a graph, and create edges to
existing vertices with a probability based on vertex i's degree ki: p(i)
= kαi + 1∑︁n j=0(kαj + 1) The α parameter allows us to configure the
model between a sublinear preference for high-degree nodes (0
\textless{} α \textless{} 1), a linear preference (α = 1), and a
superlinear preference (α \textgreater{} 1). The +1 in the above
equation ensures some chance of connecting to any vertex, even if it
initially has a degree of zero.

\begin{enumerate}
\def\labelenumi{(\roman{enumi})}
\tightlist
\item
  Implement the generator (4 points).
\end{enumerate}

Write a graph generating function preferentialGenerator(n, a) that
returns a graph with n vertices and α = a.

\begin{enumerate}
\def\labelenumi{(\roman{enumi})}
\setcounter{enumi}{1}
\tightlist
\item
  Degree analysis (4 points).
\end{enumerate}

Generate graphs with 1,000 nodes and alpha values of 0.1, 1, and 2. Make
a plot comparing the three degree distributions. You may use histograms,
kernel-density estimates, or cumulative distribution functions. However,
your plots should be easy to interpret: select number of bins or bin
widths, log scale, or other attributes as needed to produce a
high-quality figure. We will now pivot to the random walk model of
network generation. In this model, every time we add a vertex i we add
two edges. First, an edge to an existing random vertex j. Then, with
probability p, we add an edge from i to a random neighbor of j.
Otherwise we instead add an edge to another random vertex in the graph.

\begin{enumerate}
\def\labelenumi{(\roman{enumi})}
\setcounter{enumi}{2}
\tightlist
\item
  Implement the generator (4 points).
\end{enumerate}

Write a graph generating function randomWalkGenerator(n, p) that returns
a graph with n vertices generated by random walks with probability p.

\begin{enumerate}
\def\labelenumi{(\roman{enumi})}
\setcounter{enumi}{3}
\tightlist
\item
  Clustering analysis (4 points).
\end{enumerate}

Generate graphs with 1,000 nodes and p values of 0, 0.1, 0.5, and 0.9.
Make a plot comparing the four clustering distributions. As in part
(ii), make these plots easily interpretable. And finally\ldots{}

\begin{enumerate}
\def\labelenumi{(\alph{enumi})}
\setcounter{enumi}{21}
\tightlist
\item
  Interpretation (4 points).
\end{enumerate}

Compare the graphs created by the two network growth models. What kind
of structures do they produce? In what circumstances might we want to
use preferentialGenerator as a null model, and when might the random
walk model be a more appropriate null assumption?
\end{problemstmt}

    \begin{tcolorbox}[breakable, size=fbox, boxrule=1pt, pad at break*=1mm,colback=cellbackground, colframe=cellborder]
\prompt{In}{incolor}{ }{\boxspacing}
\begin{Verbatim}[commandchars=\\\{\}]
\PY{k}{def}\PY{+w}{ }\PY{n+nf}{network\PYZus{}growth\PYZus{}model}\PY{p}{(}\PY{n}{N}\PY{p}{,}\PY{n}{alpha}\PY{p}{)}\PY{p}{:}
\PY{+w}{    }\PY{l+s+sd}{\PYZdq{}\PYZdq{}\PYZdq{}}
\PY{l+s+sd}{    Generates a Network Growth Model using preferential attachment.}

\PY{l+s+sd}{    Parameters}
\PY{l+s+sd}{    \PYZhy{}\PYZhy{}\PYZhy{}\PYZhy{}\PYZhy{}\PYZhy{}\PYZhy{}\PYZhy{}\PYZhy{}\PYZhy{}}
\PY{l+s+sd}{    N : int}
\PY{l+s+sd}{        The number of vertices of the graph to generate.}
\PY{l+s+sd}{    }
\PY{l+s+sd}{    alpha : float}
\PY{l+s+sd}{        The average number of edges to make for each node added to the graph}

\PY{l+s+sd}{    Returns}
\PY{l+s+sd}{    \PYZhy{}\PYZhy{}\PYZhy{}\PYZhy{}\PYZhy{}\PYZhy{}\PYZhy{}}
\PY{l+s+sd}{    G : networkx.Graph}
\PY{l+s+sd}{        A graph generated from the given number of nodes and connections per node added (alpha). }
\PY{l+s+sd}{        }

\PY{l+s+sd}{    Notes}
\PY{l+s+sd}{    \PYZhy{}\PYZhy{}\PYZhy{}\PYZhy{}\PYZhy{}}
\PY{l+s+sd}{    \PYZhy{} }
\PY{l+s+sd}{    Time Complexity}
\PY{l+s+sd}{    \PYZhy{}\PYZhy{}\PYZhy{}\PYZhy{}\PYZhy{}\PYZhy{}\PYZhy{}\PYZhy{}\PYZhy{}\PYZhy{}\PYZhy{}\PYZhy{}\PYZhy{}\PYZhy{}\PYZhy{}}
\PY{l+s+sd}{    The time complexity is dominated by O(N**2), where N is the number of nodes we want to add to the graph,}
\PY{l+s+sd}{    each with an \PYZsq{}alpha\PYZsq{} number of connections. Where alpha is bounded by N the number of new nodes we }
\PY{l+s+sd}{    want to provide alpha connections for.}

\PY{l+s+sd}{    Example}
\PY{l+s+sd}{    \PYZhy{}\PYZhy{}\PYZhy{}\PYZhy{}\PYZhy{}\PYZhy{}\PYZhy{}}
\PY{l+s+sd}{    \PYZgt{}\PYZgt{}\PYZgt{} N = 100}
\PY{l+s+sd}{    \PYZgt{}\PYZgt{}\PYZgt{} alpha = 3}
\PY{l+s+sd}{    \PYZgt{}\PYZgt{}\PYZgt{} G = network\PYZus{}growth\PYZus{}model(N,alpha)}
\PY{l+s+sd}{    \PYZdq{}\PYZdq{}\PYZdq{}}
    \PY{c+c1}{\PYZsh{}Instantiate G }
    \PY{n}{G} \PY{o}{=} \PY{n}{nx}\PY{o}{.}\PY{n}{Graph}\PY{p}{(}\PY{p}{)}
    \PY{c+c1}{\PYZsh{}Instantiate at least 1 node}
    \PY{n}{G}\PY{o}{.}\PY{n}{add\PYZus{}nodes\PYZus{}from}\PY{p}{(}\PY{n+nb}{range}\PY{p}{(}\PY{l+m+mi}{1}\PY{p}{)}\PY{p}{)}

    \PY{c+c1}{\PYZsh{}Concept:}
    \PY{c+c1}{\PYZsh{}For every new node we want to add,}
        \PY{c+c1}{\PYZsh{}First if there are no connections then just add connect to what we initialized with}
    \PY{c+c1}{\PYZsh{}Find the probabilities of connecting to all the existing nodes currently}
    \PY{c+c1}{\PYZsh{}Pick \PYZsq{}alpha\PYZsq{} nodes at random based on said probabilities}
    \PY{c+c1}{\PYZsh{}Add a new node and connect it to each of the list of \PYZsq{}alpha\PYZsq{} nodes}
    \PY{c+c1}{\PYZsh{}Repeat for all \PYZsq{}N\PYZsq{} new nodes}

    \PY{k}{while} \PY{n}{N}\PY{o}{\PYZgt{}}\PY{l+m+mi}{0}\PY{p}{:}\PY{c+c1}{\PYZsh{}While we still have nodes to add...}
        \PY{n}{nodes} \PY{o}{=} \PY{n+nb}{list}\PY{p}{(}\PY{n}{G}\PY{o}{.}\PY{n}{nodes}\PY{p}{(}\PY{p}{)}\PY{p}{)}\PY{c+c1}{\PYZsh{}Get the current list of nodes...}
        \PY{n}{degs} \PY{o}{=} \PY{n}{np}\PY{o}{.}\PY{n}{array}\PY{p}{(}\PY{p}{[}\PY{p}{(}\PY{n}{G}\PY{o}{.}\PY{n}{degree}\PY{p}{[}\PY{n}{n}\PY{p}{]}\PY{p}{)} \PY{k}{for} \PY{n}{n} \PY{o+ow}{in} \PY{n}{nodes}\PY{p}{]}\PY{p}{,} \PY{n}{dtype}\PY{o}{=}\PY{n+nb}{float}\PY{p}{)}\PY{c+c1}{\PYZsh{}And, get their probabilities.}
        
        \PY{c+c1}{\PYZsh{}Translating alpha into an N value so that on average \PYZsq{}alpha\PYZsq{} connections are made even if this is a float}
        \PY{n}{m} \PY{o}{=} \PY{n+nb}{int}\PY{p}{(}\PY{n}{np}\PY{o}{.}\PY{n}{floor}\PY{p}{(}\PY{n}{alpha}\PY{p}{)}\PY{p}{)} \PY{o}{+} \PY{p}{(}\PY{n}{np}\PY{o}{.}\PY{n}{random}\PY{o}{.}\PY{n}{rand}\PY{p}{(}\PY{p}{)} \PY{o}{\PYZlt{}} \PY{n}{alpha} \PY{o}{\PYZhy{}} \PY{n}{np}\PY{o}{.}\PY{n}{floor}\PY{p}{(}\PY{n}{alpha}\PY{p}{)}\PY{p}{)}
        
        \PY{c+c1}{\PYZsh{}Edge Case: If this is the first run where we start with a Graph with no connections}
        \PY{k}{if} \PY{n}{degs}\PY{o}{.}\PY{n}{sum}\PY{p}{(}\PY{p}{)} \PY{o}{==} \PY{l+m+mi}{0}\PY{p}{:}
            \PY{n}{chosen} \PY{o}{=} \PY{n}{nodes} \PY{c+c1}{\PYZsh{}then just instantiate to the few nodes we initialized with                     }
        \PY{k}{else}\PY{p}{:}
            \PY{c+c1}{\PYZsh{}Otherwise, pick \PYZsq{}m\PYZsq{} connections to make **without replacement**,}
            \PY{c+c1}{\PYZsh{}  by the probability calculated from their degree(s)}
            \PY{n}{chosen} \PY{o}{=} \PY{n}{np}\PY{o}{.}\PY{n}{random}\PY{o}{.}\PY{n}{choice}\PY{p}{(}\PY{n}{nodes}\PY{p}{,} \PY{n}{size}\PY{o}{=}\PY{n}{m}\PY{p}{,} \PY{n}{replace}\PY{o}{=}\PY{k+kc}{False}\PY{p}{,} \PY{n}{p}\PY{o}{=}\PY{n}{degs} \PY{o}{/} \PY{n}{degs}\PY{o}{.}\PY{n}{sum}\PY{p}{(}\PY{p}{)}\PY{p}{)}
        \PY{c+c1}{\PYZsh{}Add the new node as the next id number}
        \PY{n}{new\PYZus{}id} \PY{o}{=} \PY{n+nb}{len}\PY{p}{(}\PY{n}{nodes}\PY{p}{)}
        \PY{c+c1}{\PYZsh{}And connect it to \PYZsq{}alpha\PYZsq{} chosen nodes}
        \PY{k}{for} \PY{n}{node\PYZus{}i} \PY{o+ow}{in} \PY{n}{chosen}\PY{p}{:}
            \PY{n}{G}\PY{o}{.}\PY{n}{add\PYZus{}edge}\PY{p}{(}\PY{n+nb}{int}\PY{p}{(}\PY{n}{node\PYZus{}i}\PY{p}{)}\PY{p}{,} \PY{n}{new\PYZus{}id}\PY{p}{)}
        \PY{c+c1}{\PYZsh{}Update N to reflect that we processed 1 new node!}
        \PY{n}{N}\PY{o}{\PYZhy{}}\PY{o}{=}\PY{l+m+mi}{1}

    \PY{k}{return} \PY{n}{G}
\end{Verbatim}
\end{tcolorbox}

    \begin{tcolorbox}[breakable, size=fbox, boxrule=1pt, pad at break*=1mm,colback=cellbackground, colframe=cellborder]
\prompt{In}{incolor}{14}{\boxspacing}
\begin{Verbatim}[commandchars=\\\{\}]
\PY{c+c1}{\PYZsh{}Testing}
\PY{k+kn}{import}\PY{+w}{ }\PY{n+nn}{matplotlib}\PY{n+nn}{.}\PY{n+nn}{animation}\PY{+w}{ }\PY{k}{as}\PY{+w}{ }\PY{n+nn}{animation}
\PY{n}{N} \PY{o}{=} \PY{l+m+mi}{10}
\PY{n}{alpha} \PY{o}{=} \PY{l+m+mf}{.1}
\PY{n}{G} \PY{o}{=} \PY{n}{network\PYZus{}growth\PYZus{}model}\PY{p}{(}\PY{n}{N}\PY{p}{,}\PY{n}{alpha}\PY{p}{)}

\PY{n+nb}{print}\PY{p}{(}\PY{n}{G}\PY{o}{.}\PY{n}{degree}\PY{p}{(}\PY{p}{)}\PY{p}{)}
\PY{n}{fig}\PY{p}{,} \PY{n}{ax} \PY{o}{=} \PY{n}{plt}\PY{o}{.}\PY{n}{subplots}\PY{p}{(}\PY{l+m+mi}{1}\PY{p}{,}\PY{l+m+mi}{1}\PY{p}{,}\PY{n}{figsize}\PY{o}{=}\PY{p}{(}\PY{l+m+mi}{7}\PY{p}{,}\PY{l+m+mi}{5}\PY{p}{)}\PY{p}{,}\PY{n}{dpi}\PY{o}{=}\PY{l+m+mi}{100}\PY{p}{)}

\PY{n}{visual} \PY{o}{=} \PY{n}{nx}\PY{o}{.}\PY{n}{draw}\PY{p}{(}\PY{n}{G}\PY{p}{,} \PY{n}{pos}\PY{o}{=}\PY{n}{nx}\PY{o}{.}\PY{n}{spring\PYZus{}layout}\PY{p}{(}\PY{n}{G}\PY{p}{)}\PY{p}{,} \PY{n}{node\PYZus{}size}\PY{o}{=}\PY{l+m+mi}{150}\PY{p}{,} \PY{n}{edgecolors}\PY{o}{=}\PY{l+s+s1}{\PYZsq{}}\PY{l+s+s1}{.2}\PY{l+s+s1}{\PYZsq{}}\PY{p}{,} \PY{n}{ax}\PY{o}{=}\PY{n}{ax}\PY{p}{)}
\end{Verbatim}
\end{tcolorbox}

    \begin{Verbatim}[commandchars=\\\{\}]
[(0, 1), (1, 2), (2, 1)]
    \end{Verbatim}

    \begin{center}
    \adjustimage{max size={0.9\linewidth}{0.9\paperheight}}{assignment02_files/assignment02_12_1.png}
    \end{center}
    { \hspace*{\fill} \\}
    
    \begin{tcolorbox}[breakable, size=fbox, boxrule=1pt, pad at break*=1mm,colback=cellbackground, colframe=cellborder]
\prompt{In}{incolor}{29}{\boxspacing}
\begin{Verbatim}[commandchars=\\\{\}]
\PY{n}{N} \PY{o}{=} \PY{l+m+mi}{1000}
\PY{n}{alpha1} \PY{o}{=} \PY{l+m+mf}{.1}
\PY{n}{alpha2} \PY{o}{=} \PY{l+m+mi}{1}
\PY{n}{alpha3} \PY{o}{=} \PY{l+m+mi}{2}

\PY{n}{G1} \PY{o}{=} \PY{p}{(}\PY{n}{network\PYZus{}growth\PYZus{}model}\PY{p}{(}\PY{n}{N}\PY{p}{,}\PY{n}{alpha1}\PY{p}{)}\PY{p}{)}
\PY{n}{deg1} \PY{o}{=} \PY{n+nb}{list}\PY{p}{(}\PY{n}{G1}\PY{o}{.}\PY{n}{degree}\PY{p}{(}\PY{p}{)}\PY{p}{)}
\PY{n}{G2} \PY{o}{=} \PY{p}{(}\PY{n}{network\PYZus{}growth\PYZus{}model}\PY{p}{(}\PY{n}{N}\PY{p}{,}\PY{n}{alpha2}\PY{p}{)}\PY{p}{)}
\PY{n}{deg2} \PY{o}{=} \PY{n+nb}{list}\PY{p}{(}\PY{n}{G2}\PY{o}{.}\PY{n}{degree}\PY{p}{(}\PY{p}{)}\PY{p}{)}
\PY{n}{G3} \PY{o}{=} \PY{p}{(}\PY{n}{network\PYZus{}growth\PYZus{}model}\PY{p}{(}\PY{n}{N}\PY{p}{,}\PY{n}{alpha3}\PY{p}{)}\PY{p}{)}
\PY{n}{deg3} \PY{o}{=} \PY{n+nb}{list}\PY{p}{(}\PY{n}{G3}\PY{o}{.}\PY{n}{degree}\PY{p}{(}\PY{p}{)}\PY{p}{)}
\end{Verbatim}
\end{tcolorbox}

    \begin{tcolorbox}[breakable, size=fbox, boxrule=1pt, pad at break*=1mm,colback=cellbackground, colframe=cellborder]
\prompt{In}{incolor}{37}{\boxspacing}
\begin{Verbatim}[commandchars=\\\{\}]
\PY{n}{fig}\PY{p}{,} \PY{n}{ax} \PY{o}{=} \PY{n}{plt}\PY{o}{.}\PY{n}{subplots}\PY{p}{(}\PY{l+m+mi}{1}\PY{p}{,}\PY{l+m+mi}{3}\PY{p}{,}\PY{n}{figsize}\PY{o}{=}\PY{p}{(}\PY{l+m+mi}{8}\PY{p}{,}\PY{l+m+mf}{3.5}\PY{p}{)}\PY{p}{,}\PY{n}{dpi}\PY{o}{=}\PY{l+m+mi}{200}\PY{p}{,}\PY{n}{sharex}\PY{o}{=}\PY{k+kc}{True}\PY{p}{,}\PY{n}{sharey}\PY{o}{=}\PY{k+kc}{True}\PY{p}{)}

\PY{n}{deg\PYZus{}bins1} \PY{o}{=} \PY{n}{np}\PY{o}{.}\PY{n}{geomspace}\PY{p}{(}\PY{n}{np}\PY{o}{.}\PY{n}{max}\PY{p}{(}\PY{p}{[}\PY{n}{np}\PY{o}{.}\PY{n}{min}\PY{p}{(}\PY{n}{deg1}\PY{p}{)}\PY{p}{,}\PY{l+m+mi}{1}\PY{p}{]}\PY{p}{)}\PY{p}{,}\PY{n}{np}\PY{o}{.}\PY{n}{max}\PY{p}{(}\PY{n}{deg1}\PY{p}{)}\PY{p}{)}
\PY{n}{ax}\PY{p}{[}\PY{l+m+mi}{0}\PY{p}{]}\PY{o}{.}\PY{n}{hist}\PY{p}{(}\PY{n}{deg1}\PY{p}{,}\PY{n}{bins}\PY{o}{=}\PY{n}{deg\PYZus{}bins1}\PY{p}{,} \PY{n}{density}\PY{o}{=}\PY{k+kc}{True}\PY{p}{,}\PY{n}{color} \PY{o}{=} \PY{l+s+s1}{\PYZsq{}}\PY{l+s+s1}{b}\PY{l+s+s1}{\PYZsq{}}\PY{p}{,}\PY{n}{histtype}\PY{o}{=}\PY{l+s+s1}{\PYZsq{}}\PY{l+s+s1}{step}\PY{l+s+s1}{\PYZsq{}}\PY{p}{)}
\PY{n}{ax}\PY{p}{[}\PY{l+m+mi}{0}\PY{p}{]}\PY{o}{.}\PY{n}{loglog}\PY{p}{(}\PY{p}{)}
\PY{n}{ax}\PY{p}{[}\PY{l+m+mi}{0}\PY{p}{]}\PY{o}{.}\PY{n}{set\PYZus{}title}\PY{p}{(}\PY{l+s+s2}{\PYZdq{}}\PY{l+s+s2}{Alpha = 0.1}\PY{l+s+s2}{\PYZdq{}}\PY{p}{)}
\PY{n}{ax}\PY{p}{[}\PY{l+m+mi}{0}\PY{p}{]}\PY{o}{.}\PY{n}{set\PYZus{}xlabel}\PY{p}{(}\PY{l+s+s2}{\PYZdq{}}\PY{l+s+s2}{Degree}\PY{l+s+s2}{\PYZdq{}}\PY{p}{)}
\PY{n}{ax}\PY{p}{[}\PY{l+m+mi}{0}\PY{p}{]}\PY{o}{.}\PY{n}{set\PYZus{}ylabel}\PY{p}{(}\PY{l+s+s2}{\PYZdq{}}\PY{l+s+s2}{Probability}\PY{l+s+s2}{\PYZdq{}}\PY{p}{)}
\PY{n}{ax}\PY{p}{[}\PY{l+m+mi}{0}\PY{p}{]}\PY{o}{.}\PY{n}{grid}\PY{p}{(}\PY{k+kc}{True}\PY{p}{)}

\PY{n}{deg\PYZus{}bins2} \PY{o}{=} \PY{n}{np}\PY{o}{.}\PY{n}{geomspace}\PY{p}{(}\PY{n}{np}\PY{o}{.}\PY{n}{max}\PY{p}{(}\PY{p}{[}\PY{n}{np}\PY{o}{.}\PY{n}{min}\PY{p}{(}\PY{n}{deg2}\PY{p}{)}\PY{p}{,}\PY{l+m+mi}{1}\PY{p}{]}\PY{p}{)}\PY{p}{,}\PY{n}{np}\PY{o}{.}\PY{n}{max}\PY{p}{(}\PY{n}{deg2}\PY{p}{)}\PY{p}{)}
\PY{n}{ax}\PY{p}{[}\PY{l+m+mi}{1}\PY{p}{]}\PY{o}{.}\PY{n}{hist}\PY{p}{(}\PY{n}{deg2}\PY{p}{,}\PY{n}{bins}\PY{o}{=}\PY{n}{deg\PYZus{}bins2}\PY{p}{,} \PY{n}{density}\PY{o}{=}\PY{k+kc}{True}\PY{p}{,}\PY{n}{color} \PY{o}{=} \PY{l+s+s1}{\PYZsq{}}\PY{l+s+s1}{b}\PY{l+s+s1}{\PYZsq{}}\PY{p}{,}\PY{n}{histtype}\PY{o}{=}\PY{l+s+s1}{\PYZsq{}}\PY{l+s+s1}{step}\PY{l+s+s1}{\PYZsq{}}\PY{p}{)}
\PY{n}{ax}\PY{p}{[}\PY{l+m+mi}{1}\PY{p}{]}\PY{o}{.}\PY{n}{loglog}\PY{p}{(}\PY{p}{)}
\PY{n}{ax}\PY{p}{[}\PY{l+m+mi}{1}\PY{p}{]}\PY{o}{.}\PY{n}{set\PYZus{}title}\PY{p}{(}\PY{l+s+s2}{\PYZdq{}}\PY{l+s+s2}{Alpha = 1}\PY{l+s+s2}{\PYZdq{}}\PY{p}{)}
\PY{n}{ax}\PY{p}{[}\PY{l+m+mi}{1}\PY{p}{]}\PY{o}{.}\PY{n}{set\PYZus{}xlabel}\PY{p}{(}\PY{l+s+s2}{\PYZdq{}}\PY{l+s+s2}{Degree}\PY{l+s+s2}{\PYZdq{}}\PY{p}{)}
\PY{c+c1}{\PYZsh{}ax[1].set\PYZus{}ylabel(\PYZdq{}Probability\PYZdq{})}
\PY{n}{ax}\PY{p}{[}\PY{l+m+mi}{1}\PY{p}{]}\PY{o}{.}\PY{n}{grid}\PY{p}{(}\PY{k+kc}{True}\PY{p}{)}

\PY{n}{deg\PYZus{}bins3} \PY{o}{=} \PY{n}{np}\PY{o}{.}\PY{n}{geomspace}\PY{p}{(}\PY{n}{np}\PY{o}{.}\PY{n}{max}\PY{p}{(}\PY{p}{[}\PY{n}{np}\PY{o}{.}\PY{n}{min}\PY{p}{(}\PY{n}{deg3}\PY{p}{)}\PY{p}{,}\PY{l+m+mi}{1}\PY{p}{]}\PY{p}{)}\PY{p}{,}\PY{n}{np}\PY{o}{.}\PY{n}{max}\PY{p}{(}\PY{n}{deg3}\PY{p}{)}\PY{p}{)}
\PY{n}{ax}\PY{p}{[}\PY{l+m+mi}{2}\PY{p}{]}\PY{o}{.}\PY{n}{hist}\PY{p}{(}\PY{n}{deg3}\PY{p}{,}\PY{n}{bins}\PY{o}{=}\PY{n}{deg\PYZus{}bins3}\PY{p}{,} \PY{n}{density}\PY{o}{=}\PY{k+kc}{True}\PY{p}{,}\PY{n}{color} \PY{o}{=} \PY{l+s+s1}{\PYZsq{}}\PY{l+s+s1}{b}\PY{l+s+s1}{\PYZsq{}}\PY{p}{,}\PY{n}{histtype}\PY{o}{=}\PY{l+s+s1}{\PYZsq{}}\PY{l+s+s1}{step}\PY{l+s+s1}{\PYZsq{}}\PY{p}{)}
\PY{n}{ax}\PY{p}{[}\PY{l+m+mi}{2}\PY{p}{]}\PY{o}{.}\PY{n}{loglog}\PY{p}{(}\PY{p}{)}
\PY{n}{ax}\PY{p}{[}\PY{l+m+mi}{2}\PY{p}{]}\PY{o}{.}\PY{n}{set\PYZus{}title}\PY{p}{(}\PY{l+s+s2}{\PYZdq{}}\PY{l+s+s2}{Alpha = 2}\PY{l+s+s2}{\PYZdq{}}\PY{p}{)}
\PY{n}{ax}\PY{p}{[}\PY{l+m+mi}{2}\PY{p}{]}\PY{o}{.}\PY{n}{set\PYZus{}xlabel}\PY{p}{(}\PY{l+s+s2}{\PYZdq{}}\PY{l+s+s2}{Degree}\PY{l+s+s2}{\PYZdq{}}\PY{p}{)}
\PY{c+c1}{\PYZsh{}ax[2].set\PYZus{}ylabel(\PYZdq{}Probability\PYZdq{})}
\PY{n}{ax}\PY{p}{[}\PY{l+m+mi}{2}\PY{p}{]}\PY{o}{.}\PY{n}{grid}\PY{p}{(}\PY{k+kc}{True}\PY{p}{)}
\end{Verbatim}
\end{tcolorbox}

    \begin{center}
    \adjustimage{max size={0.9\linewidth}{0.9\paperheight}}{assignment02_files/assignment02_14_0.png}
    \end{center}
    { \hspace*{\fill} \\}
    
    \begin{tcolorbox}[breakable, size=fbox, boxrule=1pt, pad at break*=1mm,colback=cellbackground, colframe=cellborder]
\prompt{In}{incolor}{ }{\boxspacing}
\begin{Verbatim}[commandchars=\\\{\}]
\PY{k}{def}\PY{+w}{ }\PY{n+nf}{random\PYZus{}walk\PYZus{}generator}\PY{p}{(}\PY{n}{N}\PY{p}{,}\PY{n}{p}\PY{p}{)}\PY{p}{:}
\PY{+w}{    }\PY{l+s+sd}{\PYZdq{}\PYZdq{}\PYZdq{}}
\PY{l+s+sd}{    Generates a Random Walk Model }

\PY{l+s+sd}{    Parameters}
\PY{l+s+sd}{    \PYZhy{}\PYZhy{}\PYZhy{}\PYZhy{}\PYZhy{}\PYZhy{}\PYZhy{}\PYZhy{}\PYZhy{}\PYZhy{}}
\PY{l+s+sd}{    N : int}
\PY{l+s+sd}{        The number of vertices of the graph to generate.}
\PY{l+s+sd}{    }
\PY{l+s+sd}{    p : float}
\PY{l+s+sd}{        The probability we make a second edge of each new vertex, a neighbor of the vertex first connected to.}

\PY{l+s+sd}{    Returns}
\PY{l+s+sd}{    \PYZhy{}\PYZhy{}\PYZhy{}\PYZhy{}\PYZhy{}\PYZhy{}\PYZhy{}}
\PY{l+s+sd}{    G : networkx.Graph}
\PY{l+s+sd}{        A graph generated from the given number of nodes and connections per node added (alpha). }
\PY{l+s+sd}{        }

\PY{l+s+sd}{    Notes}
\PY{l+s+sd}{    \PYZhy{}\PYZhy{}\PYZhy{}\PYZhy{}\PYZhy{}}
\PY{l+s+sd}{    \PYZhy{} }
\PY{l+s+sd}{    Time Complexity}
\PY{l+s+sd}{    \PYZhy{}\PYZhy{}\PYZhy{}\PYZhy{}\PYZhy{}\PYZhy{}\PYZhy{}\PYZhy{}\PYZhy{}\PYZhy{}\PYZhy{}\PYZhy{}\PYZhy{}\PYZhy{}\PYZhy{}}
\PY{l+s+sd}{    The time complexity is dominated by O(N), where N is the number of nodes we want to add to the graph,}
\PY{l+s+sd}{    each with 2 connections. }

\PY{l+s+sd}{    Example}
\PY{l+s+sd}{    \PYZhy{}\PYZhy{}\PYZhy{}\PYZhy{}\PYZhy{}\PYZhy{}\PYZhy{}}
\PY{l+s+sd}{    \PYZgt{}\PYZgt{}\PYZgt{} N = 100}
\PY{l+s+sd}{    \PYZgt{}\PYZgt{}\PYZgt{} p = .33}
\PY{l+s+sd}{    \PYZgt{}\PYZgt{}\PYZgt{} G = random\PYZus{}walk\PYZus{}generator(N,p)}
\PY{l+s+sd}{    }
\PY{l+s+sd}{    \PYZdq{}\PYZdq{}\PYZdq{}}
    \PY{c+c1}{\PYZsh{}Instantiate G }
    \PY{n}{G} \PY{o}{=} \PY{n}{nx}\PY{o}{.}\PY{n}{Graph}\PY{p}{(}\PY{p}{)}
    \PY{c+c1}{\PYZsh{}Instantiate at least 2 nodes (that are connected)}
    \PY{n}{G}\PY{o}{.}\PY{n}{add\PYZus{}edge}\PY{p}{(}\PY{l+m+mi}{0}\PY{p}{,}\PY{l+m+mi}{1}\PY{p}{)}

    \PY{c+c1}{\PYZsh{}Concept:}
    \PY{c+c1}{\PYZsh{} For every vertex we want to add (out of N),}
    \PY{c+c1}{\PYZsh{}  we add 2 edges. (i.e., m=2)}
    \PY{c+c1}{\PYZsh{} First edge is chosen at random (assuming no self loops?) n choose 1}
    \PY{c+c1}{\PYZsh{} Second edge is, if probability p, get random neighbor of j (i.e., n choose k\PYZus{}j) }
    \PY{c+c1}{\PYZsh{} else,}
    \PY{c+c1}{\PYZsh{} just repeat first edge}
    \PY{c+c1}{\PYZsh{}}

    \PY{k}{while} \PY{n}{N}\PY{o}{\PYZgt{}}\PY{l+m+mi}{0}\PY{p}{:}\PY{c+c1}{\PYZsh{}While we still have nodes to add...}
        \PY{n}{nodes} \PY{o}{=} \PY{n+nb}{list}\PY{p}{(}\PY{n}{G}\PY{o}{.}\PY{n}{nodes}\PY{p}{(}\PY{p}{)}\PY{p}{)}\PY{c+c1}{\PYZsh{}Get the current list of nodes...}
        \PY{n}{new\PYZus{}node}\PY{o}{=} \PY{n+nb}{len}\PY{p}{(}\PY{n}{nodes}\PY{p}{)}
        \PY{n}{vertex\PYZus{}j} \PY{o}{=} \PY{n}{np}\PY{o}{.}\PY{n}{random}\PY{o}{.}\PY{n}{choice}\PY{p}{(}\PY{n}{nodes}\PY{p}{,} \PY{n}{size}\PY{o}{=}\PY{l+m+mi}{1}\PY{p}{)}
        \PY{n}{nodes}\PY{o}{.}\PY{n}{remove}\PY{p}{(}\PY{o}{*}\PY{n}{vertex\PYZus{}j}\PY{p}{)}\PY{c+c1}{\PYZsh{}Prevent duplicate edges}
        \PY{k}{if} \PY{n}{np}\PY{o}{.}\PY{n}{random}\PY{o}{.}\PY{n}{rand}\PY{p}{(}\PY{p}{)} \PY{o}{\PYZgt{}} \PY{n}{p}\PY{p}{:}
            \PY{n}{G}\PY{o}{.}\PY{n}{add\PYZus{}edge}\PY{p}{(}\PY{n+nb}{int}\PY{p}{(}\PY{o}{*}\PY{n}{vertex\PYZus{}j}\PY{p}{)}\PY{p}{,}\PY{n}{new\PYZus{}node}\PY{p}{)}
            \PY{n}{vertex\PYZus{}j} \PY{o}{=} \PY{n}{np}\PY{o}{.}\PY{n}{random}\PY{o}{.}\PY{n}{choice}\PY{p}{(}\PY{n}{nodes}\PY{p}{,} \PY{n}{size}\PY{o}{=}\PY{l+m+mi}{1}\PY{p}{)}
            \PY{n}{G}\PY{o}{.}\PY{n}{add\PYZus{}edge}\PY{p}{(}\PY{n+nb}{int}\PY{p}{(}\PY{o}{*}\PY{n}{vertex\PYZus{}j}\PY{p}{)}\PY{p}{,}\PY{n}{new\PYZus{}node}\PY{p}{)}
        \PY{k}{else}\PY{p}{:}
            \PY{n}{vertex\PYZus{}j\PYZus{}neighbor} \PY{o}{=} \PY{n}{np}\PY{o}{.}\PY{n}{random}\PY{o}{.}\PY{n}{choice}\PY{p}{(}\PY{n+nb}{list}\PY{p}{(}\PY{n}{G}\PY{o}{.}\PY{n}{neighbors}\PY{p}{(}\PY{o}{*}\PY{n}{vertex\PYZus{}j}\PY{p}{)}\PY{p}{)}\PY{p}{,} \PY{n}{size}\PY{o}{=}\PY{l+m+mi}{1}\PY{p}{)}
            \PY{n}{G}\PY{o}{.}\PY{n}{add\PYZus{}edge}\PY{p}{(}\PY{n+nb}{int}\PY{p}{(}\PY{o}{*}\PY{n}{vertex\PYZus{}j}\PY{p}{)}\PY{p}{,}\PY{n}{new\PYZus{}node}\PY{p}{)}
            \PY{n}{G}\PY{o}{.}\PY{n}{add\PYZus{}edge}\PY{p}{(}\PY{o}{*}\PY{n}{vertex\PYZus{}j\PYZus{}neighbor}\PY{p}{,}\PY{n}{new\PYZus{}node}\PY{p}{)}
        \PY{n}{N}\PY{o}{\PYZhy{}}\PY{o}{=}\PY{l+m+mi}{1}
    \PY{k}{return} \PY{n}{G}
\end{Verbatim}
\end{tcolorbox}

    \begin{tcolorbox}[breakable, size=fbox, boxrule=1pt, pad at break*=1mm,colback=cellbackground, colframe=cellborder]
\prompt{In}{incolor}{175}{\boxspacing}
\begin{Verbatim}[commandchars=\\\{\}]
\PY{c+c1}{\PYZsh{}Testing}
\PY{k+kn}{import}\PY{+w}{ }\PY{n+nn}{matplotlib}\PY{n+nn}{.}\PY{n+nn}{animation}\PY{+w}{ }\PY{k}{as}\PY{+w}{ }\PY{n+nn}{animation}
\PY{n}{N} \PY{o}{=} \PY{l+m+mi}{1}
\PY{n}{p} \PY{o}{=} \PY{l+m+mf}{0.3}
\PY{n}{G} \PY{o}{=} \PY{n}{random\PYZus{}walk\PYZus{}generator}\PY{p}{(}\PY{n}{N}\PY{p}{,}\PY{n}{p}\PY{p}{)}
 
\PY{c+c1}{\PYZsh{} Expectations:}
\PY{c+c1}{\PYZsh{} whenever we roll a chance \PYZlt{} p, then we will be connected to an edge and its neighbor. i.e., v\PYZus{}ij1 will have an edge with v\PYZus{}ij2, and n\PYZus{}i will have edges with both of them}
\PY{c+c1}{\PYZsh{} We will have N+2 nodes (we need a minimum of two nodes to start connecting 2 edges for every vertex)}
\PY{c+c1}{\PYZsh{} We will have a minimum of 2 edges per node. N=1 should be a perfect triangle}

\PY{n+nb}{print}\PY{p}{(}\PY{n}{G}\PY{o}{.}\PY{n}{degree}\PY{p}{(}\PY{p}{)}\PY{p}{)}
\PY{n}{fig}\PY{p}{,} \PY{n}{ax} \PY{o}{=} \PY{n}{plt}\PY{o}{.}\PY{n}{subplots}\PY{p}{(}\PY{l+m+mi}{1}\PY{p}{,}\PY{l+m+mi}{1}\PY{p}{,}\PY{n}{figsize}\PY{o}{=}\PY{p}{(}\PY{l+m+mi}{7}\PY{p}{,}\PY{l+m+mi}{5}\PY{p}{)}\PY{p}{,}\PY{n}{dpi}\PY{o}{=}\PY{l+m+mi}{100}\PY{p}{)}

\PY{n}{visual} \PY{o}{=} \PY{n}{nx}\PY{o}{.}\PY{n}{draw}\PY{p}{(}\PY{n}{G}\PY{p}{,} \PY{n}{pos}\PY{o}{=}\PY{n}{nx}\PY{o}{.}\PY{n}{spring\PYZus{}layout}\PY{p}{(}\PY{n}{G}\PY{p}{)}\PY{p}{,} \PY{n}{node\PYZus{}size}\PY{o}{=}\PY{l+m+mi}{150}\PY{p}{,} \PY{n}{edgecolors}\PY{o}{=}\PY{l+s+s1}{\PYZsq{}}\PY{l+s+s1}{.2}\PY{l+s+s1}{\PYZsq{}}\PY{p}{,} \PY{n}{ax}\PY{o}{=}\PY{n}{ax}\PY{p}{)}
\end{Verbatim}
\end{tcolorbox}

    \begin{Verbatim}[commandchars=\\\{\}]
[(0, 2), (1, 2), (2, 2)]
    \end{Verbatim}

    \begin{center}
    \adjustimage{max size={0.9\linewidth}{0.9\paperheight}}{assignment02_files/assignment02_16_1.png}
    \end{center}
    { \hspace*{\fill} \\}
    
    \begin{tcolorbox}[breakable, size=fbox, boxrule=1pt, pad at break*=1mm,colback=cellbackground, colframe=cellborder]
\prompt{In}{incolor}{176}{\boxspacing}
\begin{Verbatim}[commandchars=\\\{\}]
\PY{n}{N} \PY{o}{=} \PY{l+m+mi}{1000}
\PY{n}{p1} \PY{o}{=} \PY{l+m+mi}{0}
\PY{n}{p2} \PY{o}{=} \PY{l+m+mf}{0.1}
\PY{n}{p3} \PY{o}{=} \PY{l+m+mf}{0.5}
\PY{n}{p4} \PY{o}{=} \PY{l+m+mf}{0.9}

\PY{n}{G1} \PY{o}{=} \PY{p}{(}\PY{n}{random\PYZus{}walk\PYZus{}generator}\PY{p}{(}\PY{n}{N}\PY{p}{,}\PY{n}{p1}\PY{p}{)}\PY{p}{)}
\PY{n}{cc1} \PY{o}{=} \PY{n+nb}{list}\PY{p}{(}\PY{n+nb}{dict}\PY{p}{(}\PY{n}{nx}\PY{o}{.}\PY{n}{clustering}\PY{p}{(}\PY{n}{G1}\PY{p}{)}\PY{p}{)}\PY{o}{.}\PY{n}{values}\PY{p}{(}\PY{p}{)}\PY{p}{)}
\PY{n}{G2} \PY{o}{=} \PY{p}{(}\PY{n}{random\PYZus{}walk\PYZus{}generator}\PY{p}{(}\PY{n}{N}\PY{p}{,}\PY{n}{p2}\PY{p}{)}\PY{p}{)}
\PY{n}{cc2} \PY{o}{=} \PY{n+nb}{list}\PY{p}{(}\PY{n+nb}{dict}\PY{p}{(}\PY{n}{nx}\PY{o}{.}\PY{n}{clustering}\PY{p}{(}\PY{n}{G2}\PY{p}{)}\PY{p}{)}\PY{o}{.}\PY{n}{values}\PY{p}{(}\PY{p}{)}\PY{p}{)}
\PY{n}{G3} \PY{o}{=} \PY{p}{(}\PY{n}{random\PYZus{}walk\PYZus{}generator}\PY{p}{(}\PY{n}{N}\PY{p}{,}\PY{n}{p3}\PY{p}{)}\PY{p}{)}
\PY{n}{cc3} \PY{o}{=} \PY{n+nb}{list}\PY{p}{(}\PY{n+nb}{dict}\PY{p}{(}\PY{n}{nx}\PY{o}{.}\PY{n}{clustering}\PY{p}{(}\PY{n}{G3}\PY{p}{)}\PY{p}{)}\PY{o}{.}\PY{n}{values}\PY{p}{(}\PY{p}{)}\PY{p}{)}
\PY{n}{G4} \PY{o}{=} \PY{p}{(}\PY{n}{random\PYZus{}walk\PYZus{}generator}\PY{p}{(}\PY{n}{N}\PY{p}{,}\PY{n}{p4}\PY{p}{)}\PY{p}{)}
\PY{n}{cc4} \PY{o}{=} \PY{n+nb}{list}\PY{p}{(}\PY{n+nb}{dict}\PY{p}{(}\PY{n}{nx}\PY{o}{.}\PY{n}{clustering}\PY{p}{(}\PY{n}{G4}\PY{p}{)}\PY{p}{)}\PY{o}{.}\PY{n}{values}\PY{p}{(}\PY{p}{)}\PY{p}{)}

\end{Verbatim}
\end{tcolorbox}

    \begin{tcolorbox}[breakable, size=fbox, boxrule=1pt, pad at break*=1mm,colback=cellbackground, colframe=cellborder]
\prompt{In}{incolor}{ }{\boxspacing}
\begin{Verbatim}[commandchars=\\\{\}]
\PY{n}{fig}\PY{p}{,} \PY{n}{ax} \PY{o}{=} \PY{n}{plt}\PY{o}{.}\PY{n}{subplots}\PY{p}{(}\PY{l+m+mi}{1}\PY{p}{,}\PY{l+m+mi}{4}\PY{p}{,}\PY{n}{figsize}\PY{o}{=}\PY{p}{(}\PY{l+m+mi}{8}\PY{p}{,}\PY{l+m+mf}{3.5}\PY{p}{)}\PY{p}{,}\PY{n}{dpi}\PY{o}{=}\PY{l+m+mi}{200}\PY{p}{,}\PY{n}{sharex}\PY{o}{=}\PY{k+kc}{True}\PY{p}{,}\PY{n}{sharey}\PY{o}{=}\PY{k+kc}{True}\PY{p}{)}

\PY{c+c1}{\PYZsh{}Since there are cases where the minimum CC is 0, }
\PY{c+c1}{\PYZsh{} I\PYZsq{}m starting the plot with the second\PYZhy{}most minimum}

\PY{c+c1}{\PYZsh{}Gathering 2nd most mins}
\PY{n}{cc1}\PY{o}{=}\PY{n}{np}\PY{o}{.}\PY{n}{asarray}\PY{p}{(}\PY{n}{cc1}\PY{p}{)}
\PY{n}{cc1\PYZus{}min2}\PY{o}{=} \PY{n}{np}\PY{o}{.}\PY{n}{amin}\PY{p}{(}\PY{n}{cc1}\PY{p}{[}\PY{n}{cc1} \PY{o}{!=} \PY{n}{np}\PY{o}{.}\PY{n}{amin}\PY{p}{(}\PY{n}{cc1}\PY{p}{)}\PY{p}{]}\PY{p}{)}
\PY{n}{cc2}\PY{o}{=}\PY{n}{np}\PY{o}{.}\PY{n}{asarray}\PY{p}{(}\PY{n}{cc2}\PY{p}{)}
\PY{n}{cc2\PYZus{}min2}\PY{o}{=} \PY{n}{np}\PY{o}{.}\PY{n}{amin}\PY{p}{(}\PY{n}{cc2}\PY{p}{[}\PY{n}{cc2} \PY{o}{!=} \PY{n}{np}\PY{o}{.}\PY{n}{amin}\PY{p}{(}\PY{n}{cc2}\PY{p}{)}\PY{p}{]}\PY{p}{)}
\PY{n}{cc3}\PY{o}{=}\PY{n}{np}\PY{o}{.}\PY{n}{asarray}\PY{p}{(}\PY{n}{cc3}\PY{p}{)}
\PY{n}{cc3\PYZus{}min2}\PY{o}{=} \PY{n}{np}\PY{o}{.}\PY{n}{amin}\PY{p}{(}\PY{n}{cc3}\PY{p}{[}\PY{n}{cc3} \PY{o}{!=} \PY{n}{np}\PY{o}{.}\PY{n}{amin}\PY{p}{(}\PY{n}{cc3}\PY{p}{)}\PY{p}{]}\PY{p}{)}
\PY{n}{cc4}\PY{o}{=}\PY{n}{np}\PY{o}{.}\PY{n}{asarray}\PY{p}{(}\PY{n}{cc4}\PY{p}{)}
\PY{n}{cc4\PYZus{}min2}\PY{o}{=} \PY{n}{np}\PY{o}{.}\PY{n}{amin}\PY{p}{(}\PY{n}{cc4}\PY{p}{[}\PY{n}{cc4} \PY{o}{!=} \PY{n}{np}\PY{o}{.}\PY{n}{amin}\PY{p}{(}\PY{n}{cc4}\PY{p}{)}\PY{p}{]}\PY{p}{)}

\PY{n}{cc\PYZus{}bins1} \PY{o}{=} \PY{n}{np}\PY{o}{.}\PY{n}{geomspace}\PY{p}{(}\PY{n}{cc1\PYZus{}min2}\PY{p}{,}\PY{n}{np}\PY{o}{.}\PY{n}{max}\PY{p}{(}\PY{n}{cc1}\PY{p}{)}\PY{p}{)}
\PY{n}{ax}\PY{p}{[}\PY{l+m+mi}{0}\PY{p}{]}\PY{o}{.}\PY{n}{hist}\PY{p}{(}\PY{n}{cc1}\PY{p}{,}\PY{n}{bins}\PY{o}{=}\PY{n}{cc\PYZus{}bins1}\PY{p}{,} \PY{n}{density}\PY{o}{=}\PY{k+kc}{True}\PY{p}{,}\PY{n}{color} \PY{o}{=} \PY{l+s+s1}{\PYZsq{}}\PY{l+s+s1}{b}\PY{l+s+s1}{\PYZsq{}}\PY{p}{,}\PY{n}{histtype}\PY{o}{=}\PY{l+s+s1}{\PYZsq{}}\PY{l+s+s1}{step}\PY{l+s+s1}{\PYZsq{}}\PY{p}{)}
\PY{n}{ax}\PY{p}{[}\PY{l+m+mi}{0}\PY{p}{]}\PY{o}{.}\PY{n}{loglog}\PY{p}{(}\PY{p}{)}
\PY{n}{ax}\PY{p}{[}\PY{l+m+mi}{0}\PY{p}{]}\PY{o}{.}\PY{n}{set\PYZus{}title}\PY{p}{(}\PY{l+s+s2}{\PYZdq{}}\PY{l+s+s2}{p = 0}\PY{l+s+s2}{\PYZdq{}}\PY{p}{)}
\PY{n}{ax}\PY{p}{[}\PY{l+m+mi}{0}\PY{p}{]}\PY{o}{.}\PY{n}{set\PYZus{}xlabel}\PY{p}{(}\PY{l+s+s2}{\PYZdq{}}\PY{l+s+s2}{Clustering Coefficient}\PY{l+s+s2}{\PYZdq{}}\PY{p}{)}
\PY{n}{ax}\PY{p}{[}\PY{l+m+mi}{0}\PY{p}{]}\PY{o}{.}\PY{n}{set\PYZus{}ylabel}\PY{p}{(}\PY{l+s+s2}{\PYZdq{}}\PY{l+s+s2}{Probability}\PY{l+s+s2}{\PYZdq{}}\PY{p}{)}
\PY{n}{ax}\PY{p}{[}\PY{l+m+mi}{0}\PY{p}{]}\PY{o}{.}\PY{n}{grid}\PY{p}{(}\PY{k+kc}{True}\PY{p}{)}

\PY{n}{cc\PYZus{}bins2} \PY{o}{=} \PY{n}{np}\PY{o}{.}\PY{n}{geomspace}\PY{p}{(}\PY{n}{cc2\PYZus{}min2}\PY{p}{,}\PY{n}{np}\PY{o}{.}\PY{n}{max}\PY{p}{(}\PY{n}{cc2}\PY{p}{)}\PY{p}{)}
\PY{n}{ax}\PY{p}{[}\PY{l+m+mi}{1}\PY{p}{]}\PY{o}{.}\PY{n}{hist}\PY{p}{(}\PY{n}{cc2}\PY{p}{,}\PY{n}{bins}\PY{o}{=}\PY{n}{cc\PYZus{}bins2}\PY{p}{,} \PY{n}{density}\PY{o}{=}\PY{k+kc}{True}\PY{p}{,}\PY{n}{color} \PY{o}{=} \PY{l+s+s1}{\PYZsq{}}\PY{l+s+s1}{b}\PY{l+s+s1}{\PYZsq{}}\PY{p}{,}\PY{n}{histtype}\PY{o}{=}\PY{l+s+s1}{\PYZsq{}}\PY{l+s+s1}{step}\PY{l+s+s1}{\PYZsq{}}\PY{p}{)}
\PY{n}{ax}\PY{p}{[}\PY{l+m+mi}{1}\PY{p}{]}\PY{o}{.}\PY{n}{loglog}\PY{p}{(}\PY{p}{)}
\PY{n}{ax}\PY{p}{[}\PY{l+m+mi}{1}\PY{p}{]}\PY{o}{.}\PY{n}{set\PYZus{}title}\PY{p}{(}\PY{l+s+s2}{\PYZdq{}}\PY{l+s+s2}{p = 0.1}\PY{l+s+s2}{\PYZdq{}}\PY{p}{)}
\PY{n}{ax}\PY{p}{[}\PY{l+m+mi}{1}\PY{p}{]}\PY{o}{.}\PY{n}{set\PYZus{}xlabel}\PY{p}{(}\PY{l+s+s2}{\PYZdq{}}\PY{l+s+s2}{Clustering Coefficient}\PY{l+s+s2}{\PYZdq{}}\PY{p}{)}
\PY{c+c1}{\PYZsh{}ax[1].set\PYZus{}ylabel(\PYZdq{}Probability\PYZdq{})}
\PY{n}{ax}\PY{p}{[}\PY{l+m+mi}{1}\PY{p}{]}\PY{o}{.}\PY{n}{grid}\PY{p}{(}\PY{k+kc}{True}\PY{p}{)}

\PY{n}{cc\PYZus{}bins3} \PY{o}{=} \PY{n}{np}\PY{o}{.}\PY{n}{geomspace}\PY{p}{(}\PY{n}{cc3\PYZus{}min2}\PY{p}{,}\PY{n}{np}\PY{o}{.}\PY{n}{max}\PY{p}{(}\PY{n}{cc3}\PY{p}{)}\PY{p}{)}
\PY{n}{ax}\PY{p}{[}\PY{l+m+mi}{2}\PY{p}{]}\PY{o}{.}\PY{n}{hist}\PY{p}{(}\PY{n}{cc3}\PY{p}{,}\PY{n}{bins}\PY{o}{=}\PY{n}{cc\PYZus{}bins3}\PY{p}{,} \PY{n}{density}\PY{o}{=}\PY{k+kc}{True}\PY{p}{,}\PY{n}{color} \PY{o}{=} \PY{l+s+s1}{\PYZsq{}}\PY{l+s+s1}{b}\PY{l+s+s1}{\PYZsq{}}\PY{p}{,}\PY{n}{histtype}\PY{o}{=}\PY{l+s+s1}{\PYZsq{}}\PY{l+s+s1}{step}\PY{l+s+s1}{\PYZsq{}}\PY{p}{)}
\PY{n}{ax}\PY{p}{[}\PY{l+m+mi}{2}\PY{p}{]}\PY{o}{.}\PY{n}{loglog}\PY{p}{(}\PY{p}{)}
\PY{n}{ax}\PY{p}{[}\PY{l+m+mi}{2}\PY{p}{]}\PY{o}{.}\PY{n}{set\PYZus{}title}\PY{p}{(}\PY{l+s+s2}{\PYZdq{}}\PY{l+s+s2}{p = 0.5}\PY{l+s+s2}{\PYZdq{}}\PY{p}{)}
\PY{n}{ax}\PY{p}{[}\PY{l+m+mi}{2}\PY{p}{]}\PY{o}{.}\PY{n}{set\PYZus{}xlabel}\PY{p}{(}\PY{l+s+s2}{\PYZdq{}}\PY{l+s+s2}{Clustering Coefficient}\PY{l+s+s2}{\PYZdq{}}\PY{p}{)}
\PY{c+c1}{\PYZsh{}ax[2].set\PYZus{}ylabel(\PYZdq{}Probability\PYZdq{})}
\PY{n}{ax}\PY{p}{[}\PY{l+m+mi}{2}\PY{p}{]}\PY{o}{.}\PY{n}{grid}\PY{p}{(}\PY{k+kc}{True}\PY{p}{)}

\PY{n}{cc\PYZus{}bins4} \PY{o}{=} \PY{n}{np}\PY{o}{.}\PY{n}{geomspace}\PY{p}{(}\PY{n}{cc4\PYZus{}min2}\PY{p}{,}\PY{n}{np}\PY{o}{.}\PY{n}{max}\PY{p}{(}\PY{n}{cc4}\PY{p}{)}\PY{p}{)}
\PY{n}{ax}\PY{p}{[}\PY{l+m+mi}{3}\PY{p}{]}\PY{o}{.}\PY{n}{hist}\PY{p}{(}\PY{n}{cc4}\PY{p}{,}\PY{n}{bins}\PY{o}{=}\PY{n}{cc\PYZus{}bins4}\PY{p}{,} \PY{n}{density}\PY{o}{=}\PY{k+kc}{True}\PY{p}{,}\PY{n}{color} \PY{o}{=} \PY{l+s+s1}{\PYZsq{}}\PY{l+s+s1}{b}\PY{l+s+s1}{\PYZsq{}}\PY{p}{,}\PY{n}{histtype}\PY{o}{=}\PY{l+s+s1}{\PYZsq{}}\PY{l+s+s1}{step}\PY{l+s+s1}{\PYZsq{}}\PY{p}{)}
\PY{n}{ax}\PY{p}{[}\PY{l+m+mi}{3}\PY{p}{]}\PY{o}{.}\PY{n}{loglog}\PY{p}{(}\PY{p}{)}
\PY{n}{ax}\PY{p}{[}\PY{l+m+mi}{3}\PY{p}{]}\PY{o}{.}\PY{n}{set\PYZus{}title}\PY{p}{(}\PY{l+s+s2}{\PYZdq{}}\PY{l+s+s2}{p = 0.9}\PY{l+s+s2}{\PYZdq{}}\PY{p}{)}
\PY{n}{ax}\PY{p}{[}\PY{l+m+mi}{3}\PY{p}{]}\PY{o}{.}\PY{n}{set\PYZus{}xlabel}\PY{p}{(}\PY{l+s+s2}{\PYZdq{}}\PY{l+s+s2}{Clustering Coefficient}\PY{l+s+s2}{\PYZdq{}}\PY{p}{)}
\PY{c+c1}{\PYZsh{}ax[3].set\PYZus{}ylabel(\PYZdq{}Probability\PYZdq{})}
\PY{n}{ax}\PY{p}{[}\PY{l+m+mi}{3}\PY{p}{]}\PY{o}{.}\PY{n}{grid}\PY{p}{(}\PY{k+kc}{True}\PY{p}{)}
\end{Verbatim}
\end{tcolorbox}

    \begin{center}
    \adjustimage{max size={0.9\linewidth}{0.9\paperheight}}{assignment02_files/assignment02_18_0.png}
    \end{center}
    { \hspace*{\fill} \\}
    
    Interpretation: The network-growth model produced degree distributions
that follow a power law curve (i.e., a linear shape on a log-log plot).
This means it produced networks structures where hubs emerged and tended
to be the places where new nodes made connections to. The
randomWalkGenerator on the other hand allowed us to generate various
levels of clustering coefficient distributions by picking different p
values. A preferentialGenerator may be a better suited null model
therefore when we want to simulate or model a network where we know
there are strong hubs that dominate the network. Take the internet for
example, it's a well-known fact that the most popular websites dominate
most searches and webpage requests online, and as such many new websites
may promote their own website through ads on this site (i.e., seek out a
connection to it) through for example, running ads on Google to garner
more attention, or, posting an opinion piece on a well known newpaper
website to bring attention to your own, etc. Conversely, a
randomWalkGenerator may be more helpful when you are looking at perhaps
a relatively smaller but known to be tighter knit network. Say, a circle
of friends, people, researchers, etc. Here one can leverage the
randomWalkGenerator by tuning the p-value to the level of clustering
observed in the groups of interest, to observe this level of clustering
amongst a network of these tight-knit circles
\begin{problemstmt}
    \textbf{Problem 4 (20 points)}

We are training you to conduct original research, and a crucial
component of that is a literature review. We can conduct literature
reviews for several reasons, including familiarizing ourselves with a
topic of study, a set of related methods, or the conventional wisdom of
a field. When beginning a research project we typically want to know
whether our research question has been asked before (and if so, using
what tools and with what data), what related questions have been asked,
and what gaps in understanding remain that we can try to answer. Today
we will start small. Search Google Scholar for five papers on different
community detection algorithms. Don't use the ones we've discussed
extensively in class (label propagation, Louvain, Girvan- Newman). For
each algorithm, answer the following:

\begin{enumerate}
\def\labelenumi{(\alph{enumi})}
\item
  What measure is being optimized (how do we know a given partition is
  optimal or good? Are they use modularity maximization, or something
  else?)
\item
  Algorithmically (a high level explanation in a few sentences is fine),
  how are partitions selected / sampled?
\item
  When did this technique come out? Does the journal or conference
  suggest what field this al- gorithm is used in? Is it terribly
  obscure? (We can use citation count as a proxy for obscurity here)
\item
  Is the algorithm particularly well-suited to certain kinds of graphs?
  (Does it assume sparse graphs? A core-periphery structure? Does it
  scale well / poorly? Add a little context about where you might /
  should not use this algorithm)
\end{enumerate}

I encourage you to read selectively; you want to understand how these
algorithms work and where and how they could be used, but you don't need
to understand every check for division by zero or the full literature
review of how this approach has been used in MRI analysis and how
detected neuronal communities differ from established medical
literature. Read papers for what you need from them.
\end{problemstmt}


 5 Papers on Community Detection:

\begin{enumerate}
\def\labelenumi{\arabic{enumi}.}
\tightlist
\item
  ``Detection of Overlapping Communities in Dynamical Social Networks
  \textbar{} IEEE Conference Publication \textbar{} IEEE Xplore.''
  Accessed October 30, 2026.
  https://ieeexplore.ieee.org/abstract/document/5591234.
\end{enumerate}

\begin{enumerate}
\def\labelenumi{(\alph{enumi})}
\tightlist
\item
  There are 2 values that are used to characterize communities (i.e.,
  optimized for), the mean number of second neighbors, which reflects
  the mean number of neighbors that the community is at least a path
  length or 2 or less from, and the mean number of robust second
  neighbors, which reflects the mean number of neighbors that the
  community has at least 2 paths of length or 2 or less from. (i.e.,
  that there is more than 1 path that places it on the former list).
\end{enumerate}

This is meant to capture: 
* Ability to access easily (in two steps) to
most of the community (at least as much as the others nodes of the
community) 
* Robust access to other nodes of the community (at least as
much as the others nodes of the community).

These values are optimized for adpatively to ensure the same community
can be detected whether the original whole graph is relatively sparse or
conversely relatively dense. The authors argue other algorithms tend to
produce differing results because the density of the whole graph is
often used to create what they call a ``relative'' definition of
community.

\begin{enumerate}
\def\labelenumi{(\alph{enumi})}
\setcounter{enumi}{1}
\item
  For every time-stamped set of edges, and each edge within said set,
  attempt to integrate (i.e., compare against the 2 index values) any
  nodes outside the community (i.e., taken from the edge) into the
  community. Once complete, update previously made communities. Then, If
  neither node (off the edge) is a part of a community create a new
  community. Lastly, merge similar communities. (i.e., check the
  percentage of nodes that overlap, between a pair of communities,
  against a definiable threshold)
\item
  This algorithm came out in 2010 in an IEEE conference. It is intended
  specifically to address the problem of being able to capture
  overlapping communities which can happen in multiple fields, but here
  they key-in on social networks in particular. It has 218 citations so
  it's not terribly obscure.
\item
  This algorithm is particularly well-suited to ensure results that are
  agnostic of the graph density. It allows you to define a minimial
  community as a kind of pre-defined pattern. In other-words defining
  the smallest community you'd like to detect!
\end{enumerate}

\begin{enumerate}
\def\labelenumi{\arabic{enumi}.}
\setcounter{enumi}{1}
\tightlist
\item
  Aghaalizadeh, Saeid, Saeid Taghavi Afshord, Asgarali Bouyer, and Babak
  Anari. ``A Three-Stage Algorithm for Local Community Detection Based
  on the High Node Importance Ranking in Social Networks.'' Physica A:
  Statistical Mechanics and Its Applications 563 (February 2021):
  125420. https://doi.org/10.1016/j.physa.2020.125420.
\end{enumerate}

\begin{enumerate}
\def\labelenumi{(\alph{enumi})}
\tightlist
\item
  The main measure used in LCDR, is local community detection via high
  importance nodes ranking (RDC). The RDC index for the node α indicates
  its local importance based on the degree of the neighboring nodes.
  This is calculated from thinking about important nodes as those that
  are producers of information, that can be measured via Kirchoff's law
  (i.e., analagous to currents through resistors in a circuit).
\item
\end{enumerate}

\begin{enumerate}
\def\labelenumi{\arabic{enumi}.}
\tightlist
\item
  ``Select a number of nodes as the important nodes based on the
  proposed ranking and initial community formation.''
\item
  ``Expand the initial communities by allocating the unselected nodes
  based on the proposed proximity ranking with important nodes.''
\item
  ``Merge initial communities based on the proposed relation and extract
  the final communities.''
\end{enumerate}

\begin{enumerate}
\def\labelenumi{(\alph{enumi})}
\setcounter{enumi}{2}
\tightlist
\item
  This algorithm came out in 2020 through a Physica journal. It is
  intended to produce better stability, accuracy and performance and
  scalability than other local/fast algorthms. (i.e., O(n) time). It's
  somewhat obcure at 55 citations, though it is more recent then the
  previous paper
\item
  The algorithm is intended for use in graphs when there is no prior
  knowledge beforehand and speed is desired. It can significantly
  improve accuracy relative to other local, and fast algorithms; when
  compared on simulated and real-world networks.
\end{enumerate}

3.ACM Other Conferences. ``Efficient Community Detection in Large
Networks Using Content and Links \textbar{} Proceedings of the 22nd
International Conference on World Wide Web.'' Accessed October 30, 2026.
https://dl.acm.org/doi/10.1145/2488388.2488483.

\begin{enumerate}
\def\labelenumi{(\alph{enumi})}
\tightlist
\item
  This algorithm intended to address the problem that many networks can
  actually contain noise presented as link structure, and that by
  synthesizing link structure with content, a stronger community signal
  can be better produced. In this sense they use both an estimate on
  whether a given link is at high probability to reside in a given
  community, and that it's content is similar to other nodes in the
  community as determined via cosine similarity or Jaccard coefficient.
\item
  The idea is to: ''
\end{enumerate}

\begin{enumerate}
\def\labelenumi{\arabic{enumi}.}
\tightlist
\item
  Create content edges.
\item
  Sample the union of content edges and topological edges with bias,
  retaining only edges that are relevant in local neighborhoods.
\item
  Partition the simplified graph into clusters.'' Where content edges are
  edges that can be draw based on content similiarity where as
  topological edges are those drawn more conventionally (e.g., whether
  the two nodes had contact)
\end{enumerate}

\begin{enumerate}
\def\labelenumi{(\alph{enumi})}
\setcounter{enumi}{2}
\tightlist
\item
  This technique came out in 2013 as a part of a proceeding for the
  international conference on the world wide web. It has 248 citationsso
  it's not terribly obscure, and is intended specifically to address
  noise: i.e., false positives and negatives that are common in online
  social networks and information networks.
\item
  It's particularly well-suited to real-world networks where the nodes
  reflect social media or webpages that hold inherent content or
  hyperlinks that can be used to compare one another for similarity.
\end{enumerate}

\begin{enumerate}
\def\labelenumi{\arabic{enumi}.}
\setcounter{enumi}{3}
\tightlist
\item
  Airoldi, Edoardo M., David M. Blei, Stephen E. Fienberg, and Eric P.
  Xing. ``Mixed Membership Stochastic Blockmodels.'' arXiv:0705.4485.
  Preprint, arXiv, May 30, 2007.
  https://doi.org/10.48550/arXiv.0705.4485.
\end{enumerate}

\begin{enumerate}
\def\labelenumi{(\alph{enumi})}
\tightlist
\item
  This algorithm was made to address the inherent interdependent nature
  of relational data in real-world networks that runs counter to
  standard assumptions of, for example, communities being purely
  independent. It optimizes for a kind of membership probability vector,
  where each state of the vector represents a probability that said node
  may be part of each given community.
\item
  The high-level overview is: for each node, first draw a mixed
  membership vector. Then for each pair of nodes, and for each outgoing
  edge draw a probability of membership for the holder of said outgoing
  edge (i.e., the initiator), and for each incoming edge draw a
  probability of membership for the holder of the incoming edge (i.e.,
  the receiver). From these two values sample their interaction, as
  collapsed via Bernoulli's to determine if they share membership or
  not.
\item
  This technique came out in 2007 in a NeurIPS Proceeding Advances in
  NIPS. It has 2959 citations making it quite popular!
\item
  It's particularly well-suited for networks where membership in
  communities is known interdependent (i.e., interelated) and
  establishing multiple-communities per node is desired.
\end{enumerate}

\begin{enumerate}
\def\labelenumi{\arabic{enumi}.}
\setcounter{enumi}{4}
\tightlist
\item
  Zong, Xulin, Xiangqian Ding, and Zhen Zhou. ``Health Level
  Classification by Fusing Medical Evaluation from Multiple Social
  Networks.'' Future Generation Computer Systems 114 (January 2021):
  574--80. https://doi.org/10.1016/j.future.2020.08.008.
\end{enumerate}

\begin{enumerate}
\def\labelenumi{(\alph{enumi})}
\tightlist
\item
  This algorithm is intended to optimize for accurate prediction of the
  health level of a user. It works by characterizing a user across
  multiple soureces to assign a vector of associated attributes to the
  user. The attributes are assigned by extracting known kinds of medical
  features from a user's content (e.g., analyzing image content or
  leveraging NLP). Partitions are known to be good by synthesizing data
  across multiple similar network representations of which the user is a
  part of. E.g., multiple websites.
\item
  Communities and partitions are assigned by constructing a multiple
  affinity graph to describe a vast set of potential users. Users
  determined to be similar across the various attributes to other users
  are classified as a part of a kind of ``medically-aware'' cluster that
  is intended to represent a kind of set of similar health-statuses.
\item
  This technique came out in 2021 as a part of a journal in Future
  Generation Computer Systems. It has only 5 citations making it pretty
  obscure!
\item
  It's best when there are multiple different related datasets that
  contain medical data which the algorithm can leverage to
  cross-reference against one another, and analyse to determined health
  status via online content of individuals.
\end{enumerate}

\begin{problemstmt}
    \textbf{Problem 5 OPTIONAL (5 bonus points)}

For every problem set this semester, there will be an opportunity to
earn extra credit. Since this is a new(-ish) course, we are eager to
receive genuine feedback on the material, pace, instruction style, etc.
of the course so far. Please provide any feedback about the lectures in
recent weeks and/or this assignment
\end{problemstmt}

    I'm still really enjoying the course! I think the coding exercises here
have been good for me to develop better software development practices
especially, I've been getting in the habit now of writing out expected
behavior of the code I write as a means to do better and more robust
testing. I think as we've entered more advanced content relative to the
beginning of the semester I've also been reading through the textbook
more carefully and taking more time to understand the mathematical
theory. As an engineering background the formulations aren't always my
forte so I really appreciate the time we take to explain them and really
understanding what's going on in each concept and command we tackle.

I also really appreciated the open-ended lecture we had recently to just
provide a forum to give status updates on projects and get feedback! I
think in future iterations of the class having another periodic update
or two like that throughout the semester could be really beneficial!
Especially early-on in the semester and perhaps before the end
deliverable.

I think as far constructive feedback, I don't have too much to add
relative to last assignment. Something I think might be nice, could be
to call-out textbook chapters alongside the lecture-topics in the
syllabus, as I think that could help me in being more prepared for
content ahead of each lecture! Personally I like the structure of
knowing, topic X is all contained in ChX, so that's something I've been
doing more lately (i.e., looking up the associated chapter in the
textbook) that I think has been really helpful for the assignments!

Overall, as I said before I'm really enjoying the course, I feel like
I'm learning a ton, and I think just wish I had more hours in the day to
experiment with it all hahaha

Hope this feedback is of help, and looking forward to the rest of the
course!
    % Add a bibliography block to the postdoc
    
    
    
\end{document}
